\BourbakiExercisesSection

¶ 1. Let A be a commutative ring, E, F two A-algebras, M a left E-module, N a left F-module and G the algebra \(\mathbf{G} = \mathbf{E} \otimes_{\mathbf{A}} \mathbf{F}\) . For every ordered pair of endomorphisms \(u \in \operatorname{End}_{\mathbf{E}}(\mathbf{M})\) , \(v \in \operatorname{End}_{\mathbf{F}}(\mathbf{N})\) , there exists one and only one G-endomorphism \(w\) of \(\mathbf{M} \otimes_{\mathbf{A}} \mathbf{N}\) such that \(w(x \otimes y) = u(x) \otimes v(y)\) for \(x \in \mathbf{M}\) , \(y \in \mathbf{N}\) and we can write \(w = \phi(u \otimes v)\) , where \(\phi\) is an A-linear mapping (called canonical) of \(\operatorname{End}_{\mathbf{E}}(\mathbf{M}) \otimes_{\mathbf{A}} \operatorname{End}_{\mathbf{F}}(\mathbf{N})\) into \(\operatorname{End}_{\mathbf{G}}(\mathbf{M} \otimes_{\mathbf{A}} \mathbf{N})\) . Suppose in what follows that A is a field.

\begin{enumerate}
\def\labelenumi{(\alph{enumi})}
\item
  Show that the canonical mapping \(\phi\) is injective (consider an element \(z = \sum_{i} u_{i} \otimes v_{i}\) , where the \(v_{i} \in \operatorname{End}_{\mathbf{F}}(\mathbf{N})\) are linearly independent over A and write \((\phi(z))(x_{\alpha} \otimes y) = 0\) for every element \(x_{\alpha}\) of a basis of M over A and all \(y \in \mathbf{N}\) ).
\item
  Suppose that \(\mathbf{M}\) is a finitely generated free E-module. Show that the mapping \(\phi\) is bijective in each of the following cases: (1) E is of finite rank over the field A; (2) N is a finitely generated F-module.
\item
  Suppose that \(\mathbf{M}\) is a free E-module and that there exist a \(y \in \mathbf{N}\) and an infinite sequence \((v_n)\) of endomorphisms of the F-module \(\mathbf{N}\) such that the vector subspace (over A) of \(\mathbf{N}\) generated by the \(v_n(y)\) is of infinite rank over A. Show that, if the mapping \(\phi\) is bijective, every basis of \(\mathbf{M}\) over \(\mathbf{E}\) is necessarily finite.
\item
  Suppose still that \(\mathbf{M}\) is a free E-module. Show that, for the G-module \(\mathbf{M} \otimes_{\mathbf{A}} \mathbf{N}\) to be faithful, it is necessary and sufficient that the F-module \(\mathbf{N}\) be faithful (consider a basis of \(\mathbf{E}\) over \(\mathbf{A}\) ).
\end{enumerate}

\begin{enumerate}
\def\labelenumi{\arabic{enumi}.}
\setcounter{enumi}{1}
\tightlist
\item
  Let K be a commutative field, E, F two fields whose centre contains
\end{enumerate}

K, M a left vector space over E and N a left vector space over F; let \(G = E \otimes_{K} F\) .

\begin{enumerate}
\def\labelenumi{(\alph{enumi})}
\item
  Let \(x \neq 0\) be an element of \(\mathbf{M}\) and \(y \neq 0\) an element of \(\mathbf{N}\) . Show that \(x \otimes y\) is free in the G-module \(\mathbf{M} \otimes_{\mathbb{K}} \mathbf{N}\) (use Exercise 1(d)) and the transitivity of \(\operatorname{End}_{\mathbf{E}}(\mathbf{M})\) (resp. \(\operatorname{End}_{\mathbf{F}}(\mathbf{N})\) ) in the set of elements of \(\mathbf{M}\) (resp. \(\mathbf{N}\) ) distinct from 0).
\item
  Show that \(\mathbf{M} \otimes_{\mathbb{K}} \mathbf{N}\) is a free G-module. (If \((m_{\alpha})\) is a basis of \(\mathbf{M}\) over \(\mathbf{E}\) and \((n_{\beta})\) a basis of \(\mathbf{N}\) over \(\mathbf{F}\) , show that the sum of the G-modules \(\mathrm{G}(m_{\alpha} \otimes n_{\beta})\) is direct, using a method analogous to that of (a).)
\end{enumerate}

\BourbakiExercisesSection

\begin{enumerate}
\def\labelenumi{\arabic{enumi}.}
\item
  Let N be the Z-module Z/4Z, M the submodule 2Z/4Z of N and \(j:M\to N\) the canonical injection. Show that \(\mathsf{T}(j):\mathsf{T}(\mathsf{M})\to\mathsf{T}(\mathsf{N})\) is not injective.
\item
  Let I and J be two finite sets, \(I_{0} = I \cup \{\alpha\}\) , \(J_{0} = J \cup \{\beta\}\) , where \(\alpha\) and \(\beta\) are two elements not belonging to I and J respectively, and suppose that \(I_{0}\) and \(J_{0}\) are disjoint. Let E be a vector space of finite rank over a commutative field K and let z be an element of \(\mathbf{T}_{J}^{\mathrm{I}}(\mathbf{E})\) . For all \(i \in I\) , let \(W_{i}'\) be the subspace of \(E^{*}\) consisting of the \(x^{*} \in E^{*}\) such that \(c_{\beta}^{i}(z \otimes x^{*}) = 0\) ( \(c_{\beta}^{i}\) being the contraction of index \(i \in I\) and index \(\beta \in J_{0}\) in \(\mathbf{T}_{J_{0}}^{\mathrm{I}}(\mathbf{E})\) ) and let \(V_{i}\) be the subspace of E orthogonal to \(W_{i}'\) . For all \(j \in J\) , let \(W_{j}\) be the subspace of E consisting of the \(x \in E\) such that \(c_{j}^{\alpha}(z \otimes x) = 0\) ( \(c_{j}^{\alpha}\) being the contraction of index \(\alpha \in I_{0}\) and index \(j \in J\) in \(T_{J}^{I_{0}}(E)\) ) and let \(V_{j}'\) be the subspace of \(E^{*}\) orthogonal to \(W_{j}\) . Show that z belongs to the tensor product \(\left(\bigotimes_{i \in I} V_{i}\right) \otimes \left(\bigotimes_{j \in J} V_{j}'\right)\) canonically identified with a subspace of \(T_{J}^{I}(E)\) . Moreover, if \((\mathbf{U}_{i})_{i \in I}\) is a family of subspaces of E, \((\mathbf{U}_{j}')_{j \in J}\) a family of subspaces of \(E^{*}\) such that z belongs to the tensor product \(\left(\bigotimes_{i \in I} U_{i}\right) \otimes \left(\bigotimes_{j \in J} U_{j}'\right)\) , identified with a subspace of \(T_{J}^{I}(E)\) , then \(V_{i} \subset U_{i}\) and \(V_{j}' \subset U_{j}'\) for all \(i \in I\) and \(j \in J\) (cf.~II, §7, no.8, Proposition 17).
\item
  Let M be a finitely generated projective A-module. For every endomorphism u of M, let \(\tilde{u}\) denote the tensor \(\theta_{\mathbf{M}}^{-1}(u) \in \mathbf{M}^{*} \otimes_{\mathbf{A}} \mathbf{M}\) corresponding to it.
\end{enumerate}

\begin{enumerate}
\def\labelenumi{(\alph{enumi})}
\item
  For every element \(x \in \mathbf{M}\) , show that \(u(x) = c_2^1(x \otimes \tilde{u})\) , where \(x \otimes \tilde{u}\) belongs to \(\mathbf{M} \otimes \mathbf{M}^* \otimes \mathbf{M} = \mathsf{T}_{\langle 2\rangle}^{(1,3)}(\mathbf{M})\) .
\item
  Let \(u, v\) be two endomorphisms of \(M\) and \(w = u \circ v\) ; show that \(\tilde{w} = c_3^2(\tilde{v} \otimes \tilde{u})\) , where \(\tilde{v} \otimes \tilde{u}\) belongs to \(T_{(1,3)}^{(2,4)}(M)\) .
\item
  For every automorphism \(s\) of \(\mathbf{M}\) and every tensor \(z \in \mathbf{T}_{\mathbf{J}}^{\mathbf{I}}(\mathbf{M})\) , a tensor \(s, z \in \mathbf{T}_{\mathbf{J}}^{\mathbf{I}}(\mathbf{M})\) is defined by the condition that if \(z = \bigotimes_{k \in \mathbf{I} \cup \mathbf{J}} z_k\) where \(z_k \in \mathbf{M}\) for \(k \in \mathbf{I}\) and \(z_k \in \mathbf{M}^*\) for \(k \in \mathbf{J}\) , then \(s, z = \bigotimes_{k \in \mathbf{I} \cup \mathbf{J}} z_k'\) with \(z_k' = s(z_k)\) if \(k \in \mathbf{I}\) , \(z_{k}^{\prime}=^{t}s^{-1}(z_{k})\) if \(k\in J\) . Show that the group \(\mathbf{GL}(\mathbf{M})\) operates by the law \((s,z)\mapsto s.z\) on the A-module \(T_{J}^{I}(M)\) , the mappings \(z\mapsto s.z\) being automorphisms of this A-module. Moreover, for every ordered pair of indices \(i\in I\) , \(j\in J\) ,
\end{enumerate}

\[
c _ {j} ^ {i} (s. z) = s. c _ {j} ^ {i} (z).
\]

If M is projective and finitely generated, show, in the notation of Exercise 2, that \(s. \tilde{u} = (s \circ u \circ s^{-1})^{\sim}\) .

\begin{enumerate}
\def\labelenumi{\arabic{enumi}.}
\setcounter{enumi}{3}
\tightlist
\item
  Let M be a finitely generated projective A-module; for every bilinear mapping \(u\) of \(\mathbf{M} \times \mathbf{M}\) into \(\mathbf{M}\) , let \(\tilde{u}\) denote the tensor \(\theta_{\mathbf{M}}^{-1}(u)\) in \(\mathsf{T}_{(1,2)}^{(3)}(\mathbf{M})\) . Show that for the algebra structure on \(\mathbf{M}\) defined by \(u\) to be associative, it is necessary and sufficient that the tensors \(c_4^3 (\tilde{u} \otimes \tilde{u})\) and \(c_2^6 (\tilde{u} \otimes \tilde{u})\) correspond under the canonical associativity isomorphism of \(\mathsf{T}_{(1,2,3)}^{(4)}(\mathbf{M})\) onto \(\mathsf{T}_{(1,3,4)}^{(2)}(\mathbf{M})\) .
\end{enumerate}

¶ 5. Let E be a vector space over a commutative field K and let T(E) be the tensor algebra of E.

\begin{enumerate}
\def\labelenumi{(\alph{enumi})}
\item
  Show that \(\mathsf{T}(\mathbf{E})\) contains no divisors of 0 and that the only invertible elements in \(\mathsf{T}(\mathbf{E})\) are the scalars \(\neq 0\) . If \(\dim (\mathbf{E}) > 1\) , \(\mathsf{T}(\mathbf{E})\) is a non-commutative K-algebra.
\item
  Let u, v be two elements of T(E). Show that, if there exist two elements a, b of T(E) such that \(ua = vb \neq 0\) , one of the elements u, v is a right multiple of the other. (Consider first the case where u and v are homogeneous; reduce it to the case where E is finite-dimensional and write \(u = e_{1}u_{1} + \cdots + e_{n}u_{n}\) , \(v = e_{1}v_{1} + \cdots + e_{n}v_{n}\) , where \((e_{k})\) is a basis of E. Conclude by arguing by induction on the smallest degree of the homogeneous components of ua.) Deduce that, if \(\dim(\mathrm{E}) > 1\) , the ring T(E) admits no field of left fractions (I, § 9, Exercise 15).
\item
  Show that for two non-zero elements \(u, v\) of \(\mathsf{T}(\mathbf{E})\) to be permutable, it is necessary and sufficient that there exist a vector \(x \neq 0\) in \(\mathbf{E}\) such that \(u\) and \(v\) belong to the subalgebra of \(\mathsf{T}(\mathbf{E})\) generated by \(x\) (use (b)). Deduce that, if \(\dim(\mathbf{E}) > 1\) , the centre of \(\mathsf{T}(\mathbf{E})\) is equal to \(\mathbf{K}\) .
\end{enumerate}

¶ 6. Let K be a commutative field and E, F two K-algebras each with a unit element which is identified with the unit element l of K. Consider the K-algebra T(E ⊕ F), where E ⊕ F is considered as a vector space over K; E (resp. F) is identified with a vector subspace of T(E ⊕ F) under the canonical injection. Let \(\mathfrak{J}\) be the two-sided ideal of T(E ⊕ F) generated by the elements \(x_1 \otimes x_2 - (x_1 x_2)\) and \(y_1 \otimes y_2 - (y_1 y_2)\) , where \(x_i \in E\) , \(y_i \in F\) , \(i = 1, 2\) ; let E * F denote the quotient algebra T(E ⊕ F)/\(\mathfrak{J}\); it is called the free product of the K-algebras E and F; let \(\phi\) denote the canonical homomorphism of T(E ⊕ F) onto E * F and \(\alpha\) and \(\beta\) the restrictions of \(\phi\) to E and F; \(\alpha\) and \(\beta\) are K-algebra homomorphisms.

\begin{enumerate}
\def\labelenumi{(\alph{enumi})}
\item
  Show that for every ordered pair of K-homomorphisms \(u: \mathbf{E} \to \mathbf{G}\) , \(v: \mathbf{F} \to \mathbf{G}\) into a K-algebra \(\mathbf{G}\) , there exists one and only one K-homomorphism \(w: \mathbf{E} * \mathbf{F} \to \mathbf{G}\) such that \(u = w \circ \alpha\) and \(v = w \circ \beta\) (which justifies the terminology).
\item
  Let \(R_{n}\) be the vector subspace of \(\mathsf{T}(\mathbf{E}\oplus\mathbf{F})\) the sum of the \(\mathsf{T}^{k}(\mathbf{E}\oplus\mathbf{F})\) for \(0\leqslant k\leqslant n\) and write \(J_{n}=J\cap R_{n}\) . Show that, if we consider a basis of E consisting of l and a family \((e_{\lambda})_{\lambda\in\mathbb{L}}\) and a basis of F consisting of l and a family \((f_{\mu})_{\mu\in\mathbb{M}}\) , then a supplementary subspace of \(J_{n}+R_{n-1}\) is obtained in \(R_{n}\) by considering the subspace of \(R_{n}\) with basis the elements of the form \(z_{1}\otimes z_{2}\otimes\cdots\otimes z_{n}\) where, either \(z_{2j+1}\) is equal to one of the \(e_{\lambda}\) for \(2j+1\leqslant n\) and \(z_{2j}\) is equal to one of the \(f_{\mu}\) for \(2j\leqslant n\) , or \(z_{2j+1}\) is equal to one of the \(f_{\mu}\) for \(2j+1\leqslant n\) and \(z_{2j}\) is equal to one of the \(e_{\lambda}\) for \(2j\leqslant n\) (argue by induction on n).
\item
  Let \(\mathbf{P}_n = \phi (\mathbf{R}_n)\subset \mathbf{E}*\mathbf{F}\) ; then \(\mathbf{P}_n\subset \mathbf{P}_{n + 1},\mathbf{P}_0 = \mathbf{K},\mathbf{P}_h\mathbf{P}_k\subset \mathbf{P}_{h + k}\) and \(\mathbf{E}*\mathbf{F}\) is the union of the \(\mathbf{P}_n\) . Show that there is a canonical vector space isomorphism, for \(n\geqslant 1\)
\end{enumerate}

\[
\mathrm{P} _ {n} / \mathrm{P} _ {n - 1} \rightarrow \left(\mathrm{G} _ {1} \otimes \mathrm{G} _ {2} \otimes \dots \otimes \mathrm{G} _ {n}\right) \oplus \left(\mathrm{G} _ {2} \otimes \mathrm{G} _ {3} \otimes \dots \otimes \mathrm{G} _ {n + 1}\right)
\]

where we write \(G_{k} = E/K\) for k odd, \(G_{k} = F/K\) for k even. In particular the homomorphisms \(\alpha\) and \(\beta\) are injective, which allows us to identify E and F with sub-K-algebras of E * F. \(P_{n}^{\prime}\) (resp. \(P_{n}^{\prime\prime}\) ) is used to denote the inverse images in P of \(G_{1} \otimes G_{2} \otimes \cdots \otimes G_{n}\) (resp. \(G_{2} \otimes G_{3} \otimes \cdots \otimes G_{n+1}\) ); the height of an element \(z \in E * F\) , denoted by \(h(z)\) , is the smallest integer such that \(z \in P_{n}\) ; if \(h(z) = n\) and \(z \in P_{n}^{\prime}\) (resp. \(z \in P_{n}^{\prime\prime}\) ), z is called pure odd (resp. pure even). Then \(P_{n} = P_{n}^{\prime} + P_{n}^{\prime\prime}\) and \(P_{n}^{\prime} \cap P_{n}^{\prime\prime} = P_{n-1}\) for \(n \geqslant 1\) .

\begin{enumerate}
\def\labelenumi{(\alph{enumi})}
\setcounter{enumi}{3}
\item
  Show that, if \(u, v\) are two elements of \(\mathbf{E} * \mathbf{F}\) such that \(h(u) = r\) , \(h(v) = s\) , then \(h(uv) \leqslant h(u) + h(v)\) and that \(h(uv) = h(u) + h(v)\) unless \(u\) and \(v\) are pure and, either \(r\) is even and \(u, v\) are not of the same parity, or \(r\) is odd and \(u, v\) are of the same parity. (If for example \(r\) is even, \(u \in \mathbf{P}_r'\) , \(v \in \mathbf{P}_s'\) , show that \(uv \in \mathbf{P}_{r+s-1}\) is possible only if \(u \in \mathbf{P}_{r-1}\) or \(v \in \mathbf{P}_{s-1}\) .)
\item
  Suppose that E has no divisors of zero and, in the notation of (d), suppose that r is even, u pure even and v pure odd. Show that \(h(uv) = r + s - 1\) , unless there exist two invertible elements x, y of E such that \(uy \in P_{r-1}\) and \(xb \in P_{s-1}\) . (Let \((a_i)\) (resp. \((b_j)\) ) be the basis of the supplementary subspace of \(J_r + R_{r-1}\) in \(R_r\) (resp. of \(J_s + R_{s-1}\) in \(R_s\) ) considered in (b); write \(u = \sum_i \phi(a_i)x_i\) and \(v = \sum_j y_j \phi(b_j)\) , where \(x_i\) and \(y_j\) are in E, and show that, if \(uv \in P_{r+s-2}\) , then necessarily \(x_i y_j \in K\) .) Treat similarly the other cases where \(h(uv) < h(u) + h(v)\) when E and F are assumed to have no divisors of zero.
\item
  Deduce from (d) and (e) that, when E and F have no divisors of zero, nor does E * F.
\item
  Generalize the above definitions and results to any finite number of unital K-algebras.
\end{enumerate}

\begin{enumerate}
\def\labelenumi{\arabic{enumi}.}
\setcounter{enumi}{6}
\tightlist
\item
  Let E be a vector space of finite dimension n over a field K. For every integer \(r \geqslant 1\) , the symmetric group \(S_{r}\) operates linearly on the tensor product \(E^{\otimes r}\) by the action \((\sigma, z) \mapsto \sigma \cdot z\) such that
\end{enumerate}

\[
\sigma . \left(x _ {1} \otimes x _ {2} \otimes \dots \otimes x _ {r}\right) = x _ {\sigma^ {- 1} (1)} \otimes x _ {\sigma^ {- 1} (2)} \otimes \dots \otimes x _ {\sigma^ {- 1} (r)}
\]

for every sequence of \(r\) vectors \((x_i)_{1 \leqslant i \leqslant r}\) of \(\mathbf{E}\) . Let \(u_{\sigma}(z) = \sigma, z\) . Show that for every sequence \((v_i)_{1 \leqslant i \leqslant r}\) of \(r\) endomorphisms of \(\mathbf{E}\) ,

\[
\operatorname{Tr} (u _ {\sigma} \circ (v _ {1} \otimes v _ {2} \otimes \dots \otimes v _ {r})) = \operatorname{Tr} (w _ {1}) \operatorname{Tr} (w _ {2}) \dots \operatorname{Tr} (w _ {k})\tag{*}
\]

where the endomorphisms \(w_{j}\) of \(\mathbf{E}\) are defined as follows: \(\sigma^{-1}\) is decomposed into cycles, \(\sigma^{-1} = \lambda_1\lambda_2\ldots \lambda_k\) and, if \(\lambda_{j} = (i_{1}i_{2}\dots i_{h})\) , we write

\[
w _ {j} = v _ {i _ {1}} \circ v _ {i _ {2}} \circ \dots \circ v _ {i _ {h}}.
\]

(Reduce it to calculating the left hand side of (*) when each \(v_{i}\) is of the form \(x \mapsto \langle x, a_{h}^{*} \rangle a_{k}\) , where \((a_{j})_{1 \leqslant j \leqslant n}\) is a basis of E and \((a_{j}^{*})_{1 \leqslant j \leqslant n}\) the dual basis.)

¶ 8. Let A be an integral domain and K its field of fractions; for every A-module M, let τ(M) denote its torsion module (II, § 7, no. 10).

\begin{enumerate}
\def\labelenumi{(\alph{enumi})}
\item
  Show that for the algebra \(\mathsf{T}(\mathbf{M})\) to have no divisors of zero, it is necessary and sufficient that the A-module \(\mathsf{T}(\mathbf{M})\) be torsion-free (observe that \(\mathsf{T}(\mathbf{M})_{(\mathbb{K})}\) is isomorphic to \(\mathsf{T}(\mathbf{M}_{(\mathbb{K})})\) and use Exercise 5).
\item
  Let \(u: \mathbf{M} \to \mathbf{M}'\) be an injective A-module homomorphism. Show that \(\operatorname{Ker}(\mathsf{T}(u)) \subset \tau(\mathsf{T}(\mathbf{M}))\) ; if \(\mathbf{M}'\) is torsion-free, then \(\operatorname{Ker}(\mathsf{T}(u)) = \tau(\mathsf{T}(\mathbf{M}))\) (cf.~II, § 7, no. 10, Proposition 27).
\item
  Give an example of a torsion-free A-module M such that \(\tau(\mathsf{T}(\mathbf{M}))\) is non-zero (II, § 7, Exercise 31).
\end{enumerate}

¶ 9. Let A be a commutative ring, F the free associative algebra over A of the set of integers I = {[}1, n{]}, where n ≥ 2, and X \(_{i}\) (1 ≤ i ≤ n) the corresponding indeterminates, so that a canonical basis of F over A consists of the products X \(_{i_1}\) X \(_{i_2}\) \ldots{} X \(_{i_r}\) , where (i \(_{1}\) , \ldots, i \(_{r}\) ) is an arbitrary finite sequence of elements of I.

\begin{enumerate}
\def\labelenumi{(\alph{enumi})}
\item
  Let \(C_{2}\) denote the set of commutators \([X_{i}, X_{j}] = X_{i}X_{j} - X_{j}X_{i}\) for \(1 \leqslant i < j \leqslant n\) . Assuming that \(C_{m-1}\) is defined, \(C_{m}\) denotes the set of elements \([X_{i}, P] = X_{i}P - PX_{i}\) , where \(1 \leqslant i \leqslant n\) and P runs through \(C_{m-1}\) . Let U be the graded subalgebra of F generated by l and the union of the \(C_{m}\) for \(m \geqslant 2\) . Show that, for \(1 \leqslant i \leqslant n\) , \([X_{i}, P] \in U\) for all \(P \in U\) (confine it to the case where P is a product of elements of \(\bigcup_{m} C_{m}\) and argue by induction on the number of factors).
\item
  Suppose from now on that A is a field. For every integer \(m \geqslant 0\) , let \(\mathbf{U}_m\) be the vector A-space of homogeneous elements of degree \(m\) in U (so that \(\mathbf{U}_0 = \mathbf{A}\) , \(\mathbf{U}_1 = \{0\}\) ); choose a basis \(\mathbb{R}_{m,1}, \ldots, \mathbb{R}_{m,d_m}\) in each \(\mathbf{U}_m\) , consisting of products of elements of \(\bigcup_{k\leqslant m}\mathbf{C}_k\) . For every ordered pair of integers \((m,k)\) such that \(m\geqslant 2,1\leqslant k\leqslant d_m\) , let \(\mathbf{E}_{m,k}\) be the vector sub-A-space of U generated by the \(\mathbb{R}_{p,j}\) such that \(p > m\) or \(p = m\) and \(j\geqslant k\) ; show that \(\mathbf{E}_{m,k}\) is a two-sided ideal of U and that, if \(\mathbf{L}_{m,k}\) is the left ideal of F generated by \(\mathbf{E}_{m,k},\mathbf{L}_{m,k}\) is a two-sided ideal of F.
\item
  For every multiindex \(\alpha = (\alpha_{1},\ldots ,\alpha_{n})\in \mathbf{N}^{n}\) , we write
\end{enumerate}

\[
\mathrm{X} ^ {\alpha} = \mathrm{X} _ {1} ^ {\alpha_ {1}} \mathrm{X} _ {2} ^ {\alpha_ {2}} \dots \mathrm{X} _ {n} ^ {\alpha_ {n}}.
\]

Show that the elements \(X^{\alpha}R_{m,k}\) ( \(m > 0, l \leqslant k \leqslant d_{m}\) ) form a basis of F over K. (To prove that these elements form a generating system of the vector space F, prove that every element \(X_{i_{1}} \ldots X_{i_{r}}\) is a linear combination of them by induction on the degree r. To see that the \(X^{\alpha}R_{m,k}\) are linearly independent, argue by reductio ad absurdum by considering the maximum degree with respect to one of the \(X_{i}\) of the monomials \(X^{\alpha}\) which would appear in a linear relation between the \(X^{\alpha}R_{m,k}\) and proceed by induction on this maximum degree; for this, reduce it to the case where K is infinite (considering if necessary \(F \otimes_{K} K'\) for an extension field \(K'\) of K) and observe that when \(X_{i} + \lambda\) is substituted in an \(R_{m,k}\) for \(X_{i}\) , for some \(\lambda \in K'\) , the element \(R_{m,k}\) is unaltered).

\begin{enumerate}
\def\labelenumi{(\alph{enumi})}
\setcounter{enumi}{3}
\item
  Show that the two-sided ideal \(\mathfrak{J}(\mathbf{F})\) of \(\mathbf{F}\) generated by the commutators \([u,v] = uv - vu\) of elements of \(\mathbf{F}\) is the sum of the \(\mathbf{U}_m\) for \(m\geqslant 2\) ; \(\mathbf{F} / \mathfrak{J}(\mathbf{F})\) is isomorphic to \(\mathrm{A}[\mathrm{X}_1,\dots ,\mathrm{X}_n]\) and hence is an integral domain.
\item
  Show that the algebra U is not a free associative algebra by proving that the quotient \(\mathrm{U}/\Im(\mathrm{U})\) , where \(\Im(\mathrm{U})\) is generated by the commutators of elements of U, admits nilpotent elements \(\neq0\) . (Consider the image \(z'\) in \(\mathrm{U}/\Im(\mathrm{U})\) of the element \(z = [\mathbf{X}_{i}, \mathbf{X}_{j}]\) of U; \(z'^{2} \neq 0\) but \(z'^{4} = 0\) .)
\end{enumerate}

\begin{enumerate}
\def\labelenumi{\arabic{enumi}.}
\setcounter{enumi}{9}
\tightlist
\item
  Let A be a commutative ring,
\end{enumerate}

\[
\mathrm{M} ^ {\prime} \xrightarrow {u} \mathrm{M} \xrightarrow {v} \mathrm{M} ^ {\prime \prime} \longrightarrow 0
\]

an exact sequence of A-modules and \(n\) an integer \(>0\) . Let \(L = \operatorname{Coker}(T^n(u))\) so that there is an exact sequence

\[
\mathsf {T} ^ {n} (\mathbf {M} ^ {\prime}) \xrightarrow {\mathsf {T} ^ {n} (u)} \mathsf {T} ^ {n} (\mathbf {M}) \to \mathbf {L} \to 0.
\]

For every subset J of the set \(\{1, 2, \ldots, n\}\) , we write

\[
\mathrm{P} _ {J} = \mathrm{N} _ {1} \otimes \mathrm{N} _ {2} \otimes \dots \otimes \mathrm{N} _ {n}, \quad \text { where } \mathrm{N} _ {i} = \mathrm{M} ^ {\prime} \text { if } i \in \mathrm{J}, \mathrm{N} _ {i} = \mathrm{M} \text { if } i \notin \mathrm{J}
\]

\((\mathbf{P}_{\varnothing}\) is therefore equal to \(\mathsf{T}^n (\mathbf{M}))\) and

\[
\mathrm{Q} _ {\mathrm{J}} = \mathrm{N} _ {1} \otimes \mathrm{N} _ {2} \otimes \dots \otimes \mathrm{N} _ {n}, \quad \text { where } \mathrm{N} _ {i} = \mathrm{M} ^ {\prime} \text { if } i \in \mathrm{J}, \mathrm{N} _ {i} = \mathrm{M} ^ {\prime \prime} \text { if } i \notin \mathrm{J}
\]

\((Q_{\varnothing}\) is therefore equal to \(T^{n}(M^{\prime \prime})\) .) For every integer \(i\) such that \(0 \leqslant i \leqslant n\) , let \(P^{(i)}\) denote the submodule of \(T^{n}(M)\) generated by the union of the images of the \(\mathbf{P}_{\mathbf{J}}\) for all the subsets \(\mathbf{J}\) such that \(\operatorname{Card}(\mathbf{J}) = i\) and let \(\mathbf{L}^{(i)}\) be the canonical image of \(\mathbf{P}^{(i)}\) in \(\mathbf{L}\) .

\begin{enumerate}
\def\labelenumi{(\alph{enumi})}
\tightlist
\item
  Show that there exists a canonical surjective homomorphism
\end{enumerate}

\[
\bigoplus_ {\operatorname{Card} (J) = i} Q _ {J} \rightarrow L ^ {(i)} / L ^ {(i + 1)}
\]

(cf.~II, § 5, no. 6, Proposition 6).

\begin{enumerate}
\def\labelenumi{(\alph{enumi})}
\setcounter{enumi}{1}
\tightlist
\item
  If the sequence \(0 \to \mathbf{M}' \xrightarrow{u} \mathbf{M} \xrightarrow{v} \mathbf{M}'' \to 0\) is exact and splits, show that the homomorphism defined in (a) is bijective.
\end{enumerate}

\BourbakiExercisesSection

\begin{enumerate}
\def\labelenumi{\arabic{enumi}.}
\item
  If M is a monogenous A-module, then \(\mathsf{S}(\mathsf{M}) = \mathsf{T}(\mathsf{M})\) . Deduce an example of an injective homomorphism \(j: N \to M\) such that \(\mathsf{S}(j): \mathsf{S}(\mathsf{N}) \to \mathsf{S}(\mathsf{M})\) is not injective ( \(\S 5\) , Exercise 1).
\item
  Establish for symmetric algebras the properties analogous to those of Exercise 8(a) and (b) of § 5.
\item
  Let \(\mathbf{K}\) be a field, \(\mathbf{B}\) the polynomial ring \(\mathbf{K}[\mathbf{X}_1, \mathbf{X}_2]\) and \(\mathfrak{p}\) the ideal of \(\mathbf{B}\) generated by the polynomial \(\mathbf{X}_1^3 - \mathbf{X}_2^2\) .
\end{enumerate}

\begin{enumerate}
\def\labelenumi{(\alph{enumi})}
\item
  Show that p is prime and therefore A = B/p is an integral domain. Let a and b denote the canonical images of \(X_{1}\) and \(X_{2}\) in A; a and b are not invertible in A and \(a^{3} = b^{2}\) .
\item
  Let m be the (maximal) ideal of A generated by a and b. Show that the symmetric algebra \(\mathsf{S}(\mathsf{m})\) is not an integral domain. (Observe that m is identified with the quotient module \(A^{2}/N\) , where N is the submodule of \(A^{2}\) generated by \((b, -a)\) ; conclude that \(\mathsf{S}(\mathsf{m})\) is isomorphic to the quotient ring \(\mathbf{A}[\mathbf{U}, \mathbf{V}]/(b\mathbf{U} - a\mathbf{V})\) and show that in the polynomial ring \(\mathbf{A}[\mathbf{U}, \mathbf{V}]\) the principal ideal \((b\mathbf{V} - a\mathbf{U})\) is not prime, by considering the product \(b^{2}(b\mathbf{V}^{3} - \mathbf{U}^{3})\) .)
\end{enumerate}

¶ 4. Let A be a commutative ring and a an ideal of A. The Rees algebra of the ideal a is the subalgebra R(a) of the polynomial algebra A{[}T{]} in one indeterminate, consisting of the polynomials \(c_{0} + c_{1}\mathrm{T} + \cdots + c_{n}\mathrm{T}^{n}\) , where \(c_{0} \in \mathbf{A}\) and \(c_{k} \in \mathfrak{a}^{k}\) for every integer \(k \geqslant 1\) .

\begin{enumerate}
\def\labelenumi{(\alph{enumi})}
\item
  Show that there exists one and only one surjective A-algebra homomorphism \(r: \mathsf{S}(\mathfrak{a}) \to \mathsf{R}(\mathfrak{a})\) such that the restriction of \(r\) to \(\mathfrak{a}\) is the A-linear mapping \(x \to x\mathbb{T}\) .
\item
  Suppose that the ring A is an integral domain. Let \(a_1, \ldots, a_n\) be elements of A with \(a_n \neq 0\) ; consider the A-algebra homomorphism
\end{enumerate}

\[
u: \mathrm{A} [ \mathrm{X} _ {1}, \dots , \mathrm{X} _ {n} ] \rightarrow \mathrm{A} [ \mathrm{T} ]
\]

such that \(u(\mathbf{X}_i) = a_i\mathrm{T}\) . The kernel \(\mathfrak{r} = \operatorname{Ker}(u)\) is a graded ideal whose homogeneous component \(\mathfrak{r}_m\) of degree \(m\) consists of the homogeneous polynomials \(f\) of degree \(m\) such that \(f(a_1, \ldots, a_n) = 0\) . If \(\mathfrak{r}' \subset \mathfrak{r}\) is the ideal of \(\mathrm{A}[\mathrm{X}_1, \ldots, \mathrm{X}_n]\) generated by the homogeneous component \(\mathfrak{r}_1\) of \(\mathfrak{r}\) , show that the A-module \(\mathfrak{r}/\mathfrak{r}'\) is a torsion A-module. (To verify that, for all \(f \in \mathfrak{r}_m\) , there exists \(c \neq 0\) in A such that \(cf \in \mathfrak{r}'\) , argue by induction on \(m\) , by writing

\[
f \left(\mathrm{X} _ {1}, \dots , \mathrm{X} _ {n}\right) = \mathrm{X} _ {1} f _ {1} \left(\mathrm{X} _ {1}, \dots , \mathrm{X} _ {n}\right) + \mathrm{X} _ {2} f _ {2} \left(\mathrm{X} _ {2}, \dots , \mathrm{X} _ {n}\right) + \dots + \mathrm{X} _ {n} f _ {n} \left(\mathrm{X} _ {n}\right)
\]

where the \(f_{j}\) are homogeneous of degree \(m - 1\) ; consider the polynomial \(g(\mathbf{X}_1, \ldots, \mathbf{X}_n) = \mathbf{X}_1 f_1(a_1, \ldots, a_n) + \mathbf{X}_2 f_2(a_2, \ldots, a_n) + \cdots + \mathbf{X}_n f_n(a_n)\) , observe that \(g \in \mathfrak{r}_1 \subset \mathfrak{r}'\) and form \(a_n^{m-1} f - \mathbf{X}_n^{m-1} g\) .

\begin{enumerate}
\def\labelenumi{(\alph{enumi})}
\setcounter{enumi}{2}
\tightlist
\item
  Deduce from (b) that when A is an integral domain, the kernel of the homomorphism r of (a) is the torsion module of S(a). Deduce that, for S(a) to be an integral domain, it is necessary and sufficient that S(a) be a torsion-free A-module. In particular, if the ideal a is a projective A-module, S(a) is an integral domain and isomorphic to R(a).
\end{enumerate}

\begin{enumerate}
\def\labelenumi{\arabic{enumi}.}
\setcounter{enumi}{4}
\tightlist
\item
  Let E be a vector space of finite dimension n over a field K.
\end{enumerate}

\begin{enumerate}
\def\labelenumi{(\alph{enumi})}
\tightlist
\item
  Show that for every integer \(m\) the symmetric power \(S^m(E)\) and the vector space \(S_m'(E)\) of symmetric contravariant tensors of order \(m\) have the same dimension \(\binom{m + n - 1}{n - 1} = \binom{n + m - 1}{m}\) .
\end{enumerate}

*(b) If K is of characteristic p \textgreater{} 0, show that the vector space \(\mathbb{S}_{p}^{\prime\prime}(\mathbf{E})\) of symmetrized tensors of order p has dimension \(\binom{n}{p}\) (observe that if, in a tensor product \(x_{1} \otimes x_{2} \otimes \cdots \otimes x_{p}\) of p vectors of E, \(x_{i} = x_{j}\) for an ordered pair of distinct indices i, j, then the symmetrization of this product is zero; use the fact that if \((\mathrm{H}_{i})_{1 \leqslant i \leqslant r}\) is a partition of \(\{1, 2, \ldots, p\}\) into non-empty sets, comprising at least two sets, the subgroup of \(S_{p}\) leaving each of the \(H_{i}\) stable has order not divisible by p). On the other hand, the canonical image in \(\mathbb{S}^{p}(\mathbf{E})\) of the space \(\mathbb{S}_{p}^{\prime}(\mathbf{E})\) of symmetric tensors of order p is of dimension n and is generated by the images of the tensors \(x \otimes x \otimes \cdots \otimes x\) (p times), where \(x \in E\) (use the same remark); the canonical mapping of \(\mathbb{S}_{p}^{\prime}(\mathbf{E})\) into \(\mathbb{S}^{p}(\mathbf{E})\) is therefore not bijective although the two vector spaces have the same dimension.*

\BourbakiExercisesSection

\begin{enumerate}
\def\labelenumi{\arabic{enumi}.}
\tightlist
\item
  Let A be an integral domain and K its field of fractions.
\end{enumerate}

\begin{enumerate}
\def\labelenumi{(\alph{enumi})}
\item
  Show that the exterior power \(\bigwedge^{2}\mathbf{K}\) of the A-module \(\mathbf{K}\) reduces to 0.
\item
  For every A-module \(\mathbf{E} \subset \mathbf{K}\) , show that every exterior power \(\bigwedge \mathbf{E}\) is a torsion A-module for \(p \geqslant 2\) .
\item
  If E is the A-module defined in II, §7, Exercise 31, show that \(\bigwedge\) E is non-zero. Give an example of an integral domain A and an A-module E contained in the field of fractions K of A such that no exterior power \(\bigwedge^{p}\) E reduces to 0.
\end{enumerate}

\begin{enumerate}
\def\labelenumi{\arabic{enumi}.}
\setcounter{enumi}{1}
\item
  For every ideal a of a ring A and every A-module E, the exterior algebra \(\bigwedge(\mathrm{E}/\mathfrak{a}\mathrm{E})\) is isomorphic to \((\bigwedge\mathrm{E})/\mathfrak{a}(\bigwedge\mathrm{E})\) .
\item
  Give an example of a ring A with the following property: in the A-module \(E = A^{2}\) , there exists a submodule F such that, if \(j: F \to E\) is the canonical injection, \(\bigwedge^{2} j\) is identically zero even when neither \(\bigwedge^{2} E\) nor \(\bigwedge^{2} F\) reduces to 0 (II, §4, Exercise 2).
\item
  Let E be a free A-module of dimension \(n \geqslant 3\) . Show that if \(l \leqslant p \leqslant n - 1\) the A-modules \(\bigwedge^{p}E\) and \(\bigwedge^{n-p}E\) are isomorphic; but if \(2p \neq n\) , there exists no isomorphism \(\phi\) of \(\bigwedge^{p}E\) onto \(\bigwedge^{n-p}E\) such that for every automorphism u of E, \(\phi \circ \left(\bigwedge^{p}u\right) = \left(\bigwedge^{n-p}u\right) \circ \phi\) (if \((e_{i})\) is a basis of E, consider the automorphism u such that \(u(e_{i}) = e_{i} + e_{j}, y(e_{k}) = e_{k}\) for \(k \neq i\) , and give i and j all possible values).
\item
  Let K be a commutative field, E, F two vector spaces over K and u a linear mapping of E into F of rank r. Show that if \(p \leqslant r\) the rank of \(\bigwedge^{p}u\) is equal to \(\binom{r}{p}\) and that if p \textgreater{} r, \(\bigwedge^{p}u\) is identically zero (take in E a basis r vectors of which form a basis of a subspace supplementary to \(u^{-1}(0)\) and the rest a basis of \(u^{-1}(0)\) ).
\end{enumerate}

¶ 6. Let K be a commutative field and E a vector space over K.

\begin{enumerate}
\def\labelenumi{(\alph{enumi})}
\item
  For an element \(z = \sum_{p=0}^{\infty} z_p \left( \text{with } z_p \in \bigwedge^p \mathbf{E} \right)\) of the exterior algebra \(\bigwedge \mathbf{E}\) of \(\mathbf{E}\) to be invertible, it is necessary and sufficient that \(z_0 \neq 0\) (prove it first when \(\mathbf{E}\) is finite-dimensional and then pass to the general case, noting that for all \(z \in \bigwedge \mathbf{E}\) there exists a finite-dimensional subspace \(\mathbf{F}\) of \(\mathbf{E}\) such that \(z \in \bigwedge \mathbf{F}\) ).
\item
  Suppose that K is such that \(2 \neq 0\) in K. Show that, if E is infinite-dimensional or of even finite dimension, the centre of the algebra \(\Lambda E\) consists of the elements such that \(z_{p} = 0\) for all odd indices p.~If E is of odd dimension n, the centre of \(\Lambda E\) is the sum of the above subspace and \(\Lambda^{n}E\) . In all cases, the centre of \(\Lambda E\) is identical with the set of elements of \(\Lambda E\) permutable with all the invertible elements of \(\Lambda E\) .
\end{enumerate}

\begin{enumerate}
\def\labelenumi{\arabic{enumi}.}
\setcounter{enumi}{6}
\item
  For every A-module E, show that, if we write \([u, v] = u \wedge v - v \wedge u\) for any two elements of \(\bigwedge E\) , then for any three elements \(u, v, w\) of \(\bigwedge E\) , \([(u, v], w] = 0\) .
\item
  Let E be a free A-module.
\end{enumerate}

\begin{enumerate}
\def\labelenumi{(\alph{enumi})}
\tightlist
\item
  Show that in \(\mathbf{T}^n (\mathbf{E})\) , \(\mathfrak{J}_n''\) (no. 1) is a free sub-A-module equal to the kernel of the endomorphism \(z\mapsto a.z\) and admits a supplement which is a free A-
\end{enumerate}

module; \(\bigwedge E\) is therefore canonically isomorphic to \(\mathsf{A}_n^{\prime \prime}(\mathbf{E})\) . If also in A the equation \(2\xi = 0\) implies \(\xi = 0\) , then \(\mathsf{A}_n^{\prime \prime}(\mathbf{E}) = \mathsf{A}_n^{\prime}(\mathbf{E})\) .

*(b) Suppose that A is a field of characteristic p \textgreater{} 0 and that \(p \neq 2\) . Then the canonical image of \(\mathbf{A}_{p}^{\prime}(\mathbf{E}) = \mathbf{A}_{p}^{\prime\prime}(\mathbf{E})\) in \(\bigwedge^{p} E\) is zero.

\begin{enumerate}
\def\labelenumi{(\alph{enumi})}
\setcounter{enumi}{2}
\tightlist
\item
  Suppose that \(\mathbf{A}\) is a field of characteristic 2. Then \(\mathbf{A}_n^\prime (\mathbf{E}) = \mathbf{S}_n^\prime (\mathbf{E}),\) \(\mathbf{A}_n''(\mathbf{E}) = \mathbf{S}_n''(\mathbf{E})\) , but in general \(\mathbf{A}_n^\prime (\mathbf{E})\neq \mathbf{A}_n''(\mathbf{E})\cdot *\)
\end{enumerate}

\begin{enumerate}
\def\labelenumi{\arabic{enumi}.}
\setcounter{enumi}{8}
\tightlist
\item
  Let E be an A-module. The skewsymmetrization of a linear mapping g of \(\mathsf{T}^{n}(\mathbf{E})\) into an A-module F is by definition the linear mapping ag of \(\mathsf{T}^{n}(\mathbf{E})\) into F such that \(ag(z) = g(\boldsymbol{a}z)\) for all \(z \in \mathsf{T}^{n}(\mathbf{E})\) ; if \(\phi: E^{n} \to T^{n}(\mathbf{E})\) is the canonical mapping, the n-linear mapping \(ag \circ \phi\) is also called the skewsymmetrization of the n-linear mapping \(g \circ \phi\) . Deduce from Exercise 8(a) that if E is free every alternating n-linear mapping of \(E^{n}\) into F is the skewsymmetrization of an n-linear mapping of \(E^{n}\) into F.
\end{enumerate}

¶ 10. Let E be an A-module with a generating system of n elements. Show that every alternating n-linear mapping of \(E^{n}\) into an A-module F is the skewsymmetrization of an n-linear mapping of \(E^{n}\) into F. (Let \(a_{j}\) ( \(1 \leqslant j \leqslant n\) ) be the generators of E and \(\phi: A^{n} \rightarrow E\) the homomorphism such that \(\phi(e_{j}) = a_{j}\) for all j, where \((e_{j})\) is the canonical basis of \(A^{n}\) . If \(f: E^{n} \rightarrow F\) is an alternating n-linear mapping, so is \(f_{0}: (x_{1}, \ldots, x_{n}) \mapsto f(\phi(x_{1}), \ldots, \phi(x_{n}))\) ; show that if \(g_{0}\) is an n-linear mapping of \((A^{n})^{n}\) into F such that \(f_{0} = a g_{0}\) , \(g_{0}\) can also be written as \((x_{1}, \ldots, x_{n}) \mapsto g(\phi(x_{1}), \ldots, \phi(x_{n}))\) , where g is an n-linear mapping of \(E^{n}\) into F.)

¶ 11. Let K be a commutative field in which \(2 = 0\) , A the polynomial ring \(\mathrm{K}[\mathrm{X}, \mathrm{Y}, \mathrm{Z}]\) and E the A-module \(\mathrm{A}^3/\mathrm{M}\) , where M is the submodule of \(\mathrm{A}^3\) generated by the element \((\mathrm{X}, \mathrm{Y}, \mathrm{Z})\) of \(\mathrm{A}^3\) ; finally, let F be the quotient A-module of A by the ideal \(\mathrm{AX}^2 + \mathrm{AY}^2 + \mathrm{AZ}^2\) . Show that there exists an alternating bilinear mapping of \(\mathrm{E}^2\) into F, which is the skewsymmetrization of no bilinear mapping of \(\mathrm{E}^2\) into F, but that in \(\mathsf{T}^2(\mathbf{E})\) the module \(\mathfrak{S}_2''\) is equal to the kernel of the endomorphism \(z \mapsto az\) (consider \(\mathsf{T}^2(\mathbf{E})\) as a quotient module of \(\mathsf{T}^2 (\mathbf{A}^3)\) and show that the kernel of \(z\mapsto az\) is the canonical image of the module of symmetric tensors \(\mathsf{S}_2^{\prime}(\mathbf{A}^3)\) .

¶ 12. In the notation of Exercise 11, let B be the quotient ring

\[
\mathrm{A} / (\mathrm{AX} ^ {2} + \mathrm{AY} ^ {2} + \mathrm{AZ} ^ {2})
\]

and let G be the B-module E ⊗A B. Show that in T²(G), the module \(J_{2}^{\prime\prime}\) is distinct from the kernel of the endomorphism \(z \mapsto az\) (method analogous to that of Exercise 11).

\begin{enumerate}
\def\labelenumi{\arabic{enumi}.}
\setcounter{enumi}{12}
\tightlist
\item
  \begin{enumerate}
  \def\labelenumii{(\alph{enumii})}
  \tightlist
  \item
    Let M be a graded A-module of type N. Show that there exists on the algebra \(S(M)\) (resp. \(\bigwedge(M)\) ) one and only one graduation of type N under which \(S(M)\) (resp. \(\bigwedge(M)\) ) is a graded A-algebra and which induces on M the given graduation. Show that if the homogeneous elements of even degree in M are all zero (in which case the graduation on M is called odd), the graded algebra \(\bigwedge(M)\) thus obtained is alternating ( \(\S 4\) , no. 9, Definition 7).
  \end{enumerate}
\end{enumerate}

\begin{enumerate}
\def\labelenumi{(\alph{enumi})}
\setcounter{enumi}{1}
\tightlist
\item
  Given a graded A-module of type \(\mathbf{N}\) , consider the following universal mapping problem: \(\Sigma\) is the species of anticommutative graded algebra structure (§ 4, no. 9, Definition 7), the morphisms are A-homomorphisms of graded algebras (§ 3, no. 1) and the \(\alpha\) -mappings are the graded homomorphisms of degree 0 of M into an anticommutative graded A-algebra. Let \(\mathbf{M}^{-}\) (resp. \(\mathbf{M}^{+}\) ) denote the graded submodule of M generated by the homogeneous elements of odd (resp. even) degree. Show that the graded algebra \(\mathbf{G}(\mathbf{M}) = \bigwedge (\mathbf{M}^{-}) \otimes_{\mathbf{A}} \mathbf{S}(\mathbf{M}^{+})\) is anticommutative; a canonical injection \(j\) of M into this algebra is defined by writing \(j(x) = x \otimes 1\) if \(x \in \mathbf{M}^{-}, j(x) = 1 \otimes x\) if \(x \in \mathbf{M}^{+}\) . Show that the graded algebra \(\mathbf{G}(\mathbf{M})\) and the injection \(j\) constitute a solution of the above universal mapping problem.
\end{enumerate}

\[
\mathbf {G} (\mathbf {M}) = \bigwedge (\mathbf {M} ^ {-}) \otimes_ {\mathbf {A}} \mathbf {S} (\mathbf {M} ^ {+})
\]

is called the universal anticommutative algebra over the graded A-module M.

\begin{enumerate}
\def\labelenumi{(\alph{enumi})}
\setcounter{enumi}{2}
\tightlist
\item
  Prove the analogues of Propositions 2, 3 and 4 for the universal anticommutative algebra. Consider the case where M admits a basis consisting of homogeneous elements: find explicitly in this case a basis of the universal anticommutative algebra.
\end{enumerate}

\begin{enumerate}
\def\labelenumi{\arabic{enumi}.}
\setcounter{enumi}{13}
\tightlist
\item
  Let M be a projective A-module and let \(x_{1}, \ldots, x_{n}\) be linearly independent elements of M; for every subset \(H = \{i_{1}, \ldots, i_{m}\}\) of \(m \leqslant n\) elements of the interval \([1, n]\) , where \(i_{1} < i_{2} < \cdots < i_{m}\) , we write
\end{enumerate}

\[
x _ {\mathrm{H}} = x _ {i _ {1}} \wedge x _ {i _ {2}} \wedge \dots \wedge x _ {i _ {m}}.
\]

Show that when H runs through the set of subsets of \([1, n]\) with m elements, the \(x_{H}\) are linearly independent in \(\bigwedge(M)\) .

\begin{enumerate}
\def\labelenumi{\arabic{enumi}.}
\setcounter{enumi}{14}
\tightlist
\item
  Let M be an n-dimensional free A-module. Show that every generating system of M with n elements is a basis of M.
\end{enumerate}

\BourbakiExercisesSection

\begin{enumerate}
\def\labelenumi{\arabic{enumi}.}
\item
  In a matrix U of order n, if, for each index i, the column of index i is replaced by the sum of the columns of index \(\neq i\) , the determinant of the matrix obtained is equal to \((-1)^{n-1}(n-1)\det U\) . If in U, for each i, the sum of the columns of index \(\neq i\) is subtracted from the column of index i, the determinant obtained is equal to \(-(n-2)2^{n-1}\det U\) .
\item
  Let \(\Delta = \det (\alpha_{ij})\) be the determinant of a matrix of order \(n\) ; for
\end{enumerate}

\[
1 \leqslant i \leqslant n - 1
\]

and \(1 \leqslant j \leqslant n - 1\) , we write

\[
\beta_ {i j} = \left| \begin{array}{c c} \alpha_ {\mathbf {1} j} & \alpha_ {\mathbf {1}, j + \mathbf {1}} \\ \alpha_ {i + \mathbf {1}, j} & \alpha_ {i + \mathbf {1}, j + \mathbf {1}} \end{array} \right|.
\]

Show that the determinant \(\det(\beta_{ij})\) of order n - 1 is equal to

\[
\alpha_ {1 2} \alpha_ {1 3} \dots \alpha_ {1, n - 1} \Delta .
\]

\begin{enumerate}
\def\labelenumi{\arabic{enumi}.}
\setcounter{enumi}{2}
\tightlist
\item
  Show the identity
\end{enumerate}

\[
\left| \begin{array}{cccccc} x _ {1} & y _ {1} & \alpha_ {1 3} & \alpha_ {1 4} & \ldots & \alpha_ {1 n} \\ \lambda_ {1} x _ {2} & x _ {2} & y _ {2} & \alpha_ {2 4} & \ldots & \alpha_ {2 n} \\ \lambda_ {1} \lambda_ {2} x _ {3} & \lambda_ {2} x _ {3} & x _ {3} & y _ {3} & \ldots & \alpha_ {3 n} \\ \cdot & \cdot & \cdot & \cdot & \cdot & \cdot \\ \lambda_ {1} \lambda_ {2} \ldots \lambda_ {n - 1} x _ {n} & \lambda_ {2} \ldots \lambda_ {n - 1} x _ {n} & \lambda_ {3} \ldots \lambda_ {n - 1} x _ {n} & \lambda_ {4} \ldots \lambda_ {n - 1} x _ {n} & \ldots & x _ {n} \\ \end{array} \right| = (x _ {1} - \lambda_ {1} y _ {1}) (x _ {2} - \lambda_ {2} y _ {2}) \ldots (x _ {n - 1} - \lambda_ {n - 1} y _ {n - 1}) x _ {n}.
\]

Deduce the following identities:

\[
\left| \begin{array}{c c c c c} 1 & 1 & 1 & \dots & 1 \\ b _ {1} & a _ {1} & a _ {1} & \dots & a _ {1} \\ b _ {1} & b _ {2} & a _ {2} & \dots & a _ {2} \\ \cdot & \cdot & \cdot & \cdot & \cdot \\ b _ {1} & b _ {2} & b _ {3} & \dots & a _ {n} \end{array} \right| = (a _ {1} - b _ {1}) (a _ {2} - b _ {2}) \dots (a _ {n} - b _ {n})
\]

\[
\left| \begin{array}{c c c c c} a _ {1} b _ {1} & a _ {1} b _ {2} & a _ {1} b _ {3} & \ldots & a _ {1} b _ {n} \\ a _ {1} b _ {2} & a _ {2} b _ {2} & a _ {2} b _ {3} & \ldots & a _ {2} b _ {n} \\ a _ {1} b _ {3} & a _ {2} b _ {3} & a _ {3} b _ {3} & \ldots & a _ {3} b _ {n} \\ \cdot \cdot & \cdot \cdot & \cdot \cdot & \cdot \cdot & \cdot \cdot \\ a _ {1} b _ {n} & a _ {2} b _ {n} & a _ {3} b _ {n} & \ldots & a _ {n} b _ {n} \end{array} \right| = a _ {1} b _ {n} (a _ {2} b _ {1} - a _ {1} b _ {2}) \ldots (a _ {n} b _ {n - 1} - a _ {n - 1} b _ {n})
\]

\[
\left| \begin{array}{cccccc} a _ {2} a _ {3} \ldots a _ {n} & a _ {3} a _ {4} \ldots a _ {n} b _ {1} & a _ {4} \ldots a _ {n} b _ {1} b _ {2} & \ldots & b _ {1} b _ {2} b _ {3} \ldots b _ {n - 1} \\ b _ {2} b _ {3} \ldots b _ {n} & a _ {3} a _ {4} \ldots a _ {n} a _ {1} & a _ {4} \ldots a _ {n} a _ {1} b _ {2} & \ldots & a _ {1} b _ {2} b _ {3} \ldots b _ {n - 1} \\ a _ {2} b _ {3} \ldots b _ {n} & b _ {3} b _ {4} \ldots b _ {n} b _ {1} & a _ {4} \ldots a _ {n} a _ {1} a _ {2} & \ldots & a _ {1} a _ {2} b _ {3} \ldots b _ {n - 1} \\ \cdot & \cdot & \cdot & \cdot & \cdot \\ a _ {2} a _ {3} \ldots b _ {n} & a _ {3} a _ {4} \ldots b _ {n} b _ {1} & a _ {4} \ldots b _ {n} b _ {1} b _ {2} & \ldots & a _ {1} a _ {2} a _ {3} \ldots a _ {n - 1} \\ \end{array} \right| = (a _ {1} a _ {2} \ldots a _ {n} - b _ {1} b _ {2} \ldots b _ {n}) ^ {n - 1}.
\]

\[
\left| \begin{array}{c c c c c c} x & a _ {1} & a _ {2} & \ldots & a _ {n - 1} & 1 \\ a _ {1} & x & a _ {2} & \ldots & a _ {n - 1} & 1 \\ a _ {1} & a _ {2} & x & \ldots & a _ {n - 1} & 1 \\ . & . & . & . & . & . \\ a _ {1} & a _ {2} & a _ {3} & \ldots & x & 1 \\ a _ {1} & a _ {2} & a _ {3} & \ldots & a _ {n} & 1 \end{array} \right| = (x - a _ {1}) (x - a _ {2}) \ldots (x - a _ {n})
\]

\[
\begin{array}{r l} & {\left| \begin{array}{l l l l l} x & a _ {1} & a _ {2} & \ldots & a _ {n} \\ a _ {1} & x & a _ {2} & \ldots & a _ {n} \\ a _ {1} & a _ {2} & x & \ldots & a _ {n} \\ \cdot & \cdot & \cdot & \cdot & \cdot \\ a _ {1} & a _ {2} & a _ {3} & \ldots & x \end{array} \right|} \\ & {\qquad = (x + a _ {1} + a _ {2} + \dots + a _ {n}) (x - a _ {1}) (x - a _ {2}) \ldots (x - a _ {n})} \end{array}
\]

(reduce the last determinant to the preceding one).

\begin{enumerate}
\def\labelenumi{\arabic{enumi}.}
\setcounter{enumi}{3}
\tightlist
\item
  Calculate the determinant
\end{enumerate}

\[
\Delta_ {n} = \left| \begin{array}{c c c c c} a _ {1} + b _ {1} & b _ {1} & b _ {1} & \ldots & b _ {1} \\ b _ {2} & a _ {2} + b _ {2} & b _ {2} & \ldots & b _ {2} \\ \cdot & \cdot & \cdot & \cdot & \cdot \\ b _ {n} & b _ {n} & b _ {n} & \ldots & a _ {n} + b _ {n} \end{array} \right|
\]

(express \(\Delta_{n}\) in terms of \(\Delta_{n-1}\) ).

\begin{enumerate}
\def\labelenumi{\arabic{enumi}.}
\setcounter{enumi}{4}
\tightlist
\item
  If the \(a_i\) and \(b_j\) are elements of a commutative field such that \(a_i + b_j \neq 0\) for every ordered pair of indices \((i,j)\) , prove that
\end{enumerate}

\[
\det \left(\frac {1}{a _ {i} + b _ {j}}\right) = \frac {\sum_ {i <   j} (a _ {j} - a _ {i}) (b _ {j} - b _ {i})}{\sum_ {i , j} (a _ {i} + b _ {j})}
\]

(``Cauchy's determinant'').

\begin{enumerate}
\def\labelenumi{\arabic{enumi}.}
\setcounter{enumi}{5}
\tightlist
\item
  Show that, if X is a matrix with n rows and m columns, Y a matrix with p rows and n columns and Z = YX the product matrix with p rows and m columns, the minors of Z of order q are zero if n \textless{} q; if \(q \leqslant n\) , they are given by
\end{enumerate}

\[
\det (Z _ {L, H}) = \sum_ {K} \det (Y _ {L, K}) \det (X _ {K, H})
\]

where K runs through the set of subsets of \([1, n]\) with q elements (use formula (3) of §7, no. 2).

\begin{enumerate}
\def\labelenumi{\arabic{enumi}.}
\setcounter{enumi}{6}
\item
  Let \(\Delta = \det(\alpha_{ij})\) be the determinant of a matrix U of order n; for each index i ( \(1 \leqslant i \leqslant n\) ) let \(\Delta_i\) denote the determinant obtained by multiplying the element \(\alpha_{ij}\) in U by \(\beta_j\) for \(1 \leqslant j \leqslant n\) . Show that the sum of the n determinants \(\Delta_i\) ( \(1 \leqslant i \leqslant n\) ) is equal to \((\beta_1 + \beta_2 + \cdots + \beta_n)\Delta\) (expand \(\Delta_i\) by the row of index i).
\item
  Let \(\Delta = \det(\alpha_{ij})\) be the determinant of a matrix \(U\) of order \(n\) and \(\sigma\) a permutation belonging to \(\mathfrak{S}_n\) ; for each index \(i\) ( \(1 \leqslant i \leqslant n\) ) let \(\Delta_i\) be the determinant obtained by replacing the element \(\alpha_{ij}\) in \(U\) by \(\alpha_{i,\sigma(j)}\) for \(1 \leqslant j \leqslant n\) ; if \(p\) is the number of indices invariant under the permutation \(\sigma\) , show that the sum of the \(n\) determinants \(\Delta_i\) ( \(1 \leqslant i \leqslant n\) ) is equal to \(p\Delta\) (same method as in Exercise 7).
\item
  Let A be a square matrix of order n, B a submatrix of A with p rows and q columns and C the matrix obtained by multiplying in A, each of the elements belonging to B by the same scalar \(\alpha\) . Show that each term in the total expansion of det C is equal to the corresponding term in the total expansion of det A multiplied by a scalar of the form \(\alpha^{r}\) , where \(r \geqslant p + q - n\) and r depends on the term considered (perform a suitable Laplace expansion of det C).
\item
  Let \(\Gamma, \Delta\) be the determinants of two matrices \(U\) , \(V\) of order \(n\) ; let H and K be any two subsets of \([1, n]\) with \(p\) elements and \((i_k)\) (resp. \((j_k)\) ) the sequence obtained by arranging the elements of H (resp. K) in increasing order; let \(\Gamma_{H,K}\) be the determinant obtained by replacing in \(U\) the column of index \(i_k\) by the column of \(V\) of index \(j_k\) , for \(1 \leqslant k \leqslant p\) ; similarly let \(\Delta_{K,H}\) be the determinant obtained by replacing in \(V\) the column of index \(j_k\) by the column of \(U\) of index \(i_k\) , for \(1 \leqslant k \leqslant p\) . Show that, for every subset \(H\) of \([1, n]\) with \(p\) elements,
\end{enumerate}

\[
\Gamma \Delta = \sum_ {\mathbf {K}} \Gamma_ {\mathbf {H}, \mathbf {K}} \Delta_ {\mathbf {K}, \mathbf {H}}
\]

where K runs through the set of subsets of \([1, n]\) with p elements (use formulae (21) and (22) of no. 6).

\begin{enumerate}
\def\labelenumi{\arabic{enumi}.}
\setcounter{enumi}{10}
\tightlist
\item
  Let \(\Delta\) be the determinant of a square matrix \(X\) of order \(n\) and \(\Delta_p\) the determinant of the square matrix \(\bigwedge^p X\) , of order \(\binom{n}{p}\) . Show that
\end{enumerate}

\[
\Delta_ {p} \Delta_ {n - p} = \Delta_ {(p)} ^ {(n)}
\]

(use formulae (21) and (22) of no. 6).

\begin{enumerate}
\def\labelenumi{\arabic{enumi}.}
\setcounter{enumi}{11}
\tightlist
\item
  A matrix \((\alpha_{ij})\) of order \(n\) is called centrosymmetric (resp. skew-centrosymmetric) if \(\alpha_{n - i + 1, n - j + 1} = \alpha_{ij}\) (resp. \(\alpha_{n - i + 1, n - j + 1} = -\alpha_{ij}\) ) for \(1 \leqslant i \leqslant n\) , \(1 \leqslant j \leqslant n\) .
\end{enumerate}

\begin{enumerate}
\def\labelenumi{(\alph{enumi})}
\item
  Show that the determinant of a centrosymmetric matrix of even order 2p can be expressed in the form of a product of two determinants of order p and the determinant of a centrosymmetric matrix of order \(2p + 1\) in the form of a product of a determinant of order p and a determinant of order \(p + 1\) .
\item
  Show that the determinant of a skew-centrosymmetric matrix of even order \(2p\) can be expressed in the form of a product of two determinants of order \(p\) . The determinant of a skew-centrosymmetric matrix of odd order \(2p + 1\) is zero if in A the relation \(2\xi = 0\) implies \(\xi = 0\) ; otherwise it can be expressed in the form of the product of \(\alpha_{p+1,p+1}\) by two determinants of order \(p\) .
\end{enumerate}

\begin{enumerate}
\def\labelenumi{\arabic{enumi}.}
\setcounter{enumi}{12}
\tightlist
\item
  Let \(\Delta = \det(\alpha_{ij})\) be a determinant of order \(n\) and \(\Delta_{ij}\) the minor of order \(n - 1\) , the complement of \(\alpha_{ij}\) . Show that
\end{enumerate}

\[
\left| \begin{array}{c c c c c} \alpha_ {1 1} & \alpha_ {1 2} & \ldots & \alpha_ {1 n} & x _ {1} \\ \alpha_ {2 1} & \alpha_ {2 2} & \ldots & \alpha_ {2 n} & x _ {2} \\ \cdot & \cdot & \cdot & \cdot & \cdot \\ \alpha_ {n 1} & \alpha_ {n 2} & \ldots & \alpha_ {n n} & x _ {n} \\ y _ {1} & y _ {2} & \ldots & y _ {n} & z \end{array} \right| = \Delta z - \sum_ {i, j} (- 1) ^ {i + j} \Delta_ {i j} x _ {i} y _ {j}.
\]

If \(\Delta = 0\) and the \(\alpha_{ij}\) belong to a field, show that the above determinant is the product of a linear form in \(x_1, x_2, \ldots, x_n\) and a linear form in \(y_1, y_2, \ldots, y_n\) (use Exercise 5 of §5 and Exercise 8 of II, §6).

Give an example where this result fails to hold when the ring of scalars A is not a field (take A to be the ring \(\mathbf{Z} / (6)\) and \(n = 2\) ).

\begin{enumerate}
\def\labelenumi{\arabic{enumi}.}
\setcounter{enumi}{13}
\tightlist
\item
  Show the identity
\end{enumerate}

\[
\left| \begin{array}{cccccc} 0 & 1 & 1 & 1 & \ldots & 1 \\ 1 & 0 & a _ {1} + a _ {2} & a _ {1} + a _ {3} & \ldots & a _ {1} + a _ {n} \\ 1 & a _ {2} + a _ {1} & 0 & a _ {2} + a _ {3} & \ldots & a _ {2} + a _ {n} \\ 1 & a _ {3} + a _ {1} & a _ {3} + a _ {2} & 0 & \ldots & a _ {3} + a _ {n} \\ \cdot & \cdot & \cdot & \cdot & \cdot & \cdot \\ 1 & a _ {n} + a _ {1} & a _ {n} + a _ {2} & a _ {n} + a _ {3} & \ldots & 0 \\ \end{array} \right| = (- 1) ^ {n} 2 ^ {n - 1} \sum_ {i = 1} ^ {n} a _ {1} a _ {2} \ldots a _ {i - 1} a _ {i + 1} \ldots a _ {n}
\]

(use Exercise 13).

\begin{enumerate}
\def\labelenumi{\arabic{enumi}.}
\setcounter{enumi}{14}
\item
  Show that if the columns of a square matrix of order n over A are linearly independent, the rows of the matrix are also linearly independent (cf.~II, § 10, Exercise 3).
\item
  Let E, F be two free A-modules of respective dimensions m and n.~For a linear mapping \(u: E \rightarrow F\) to be injective, it is necessary and sufficient that \(m \leqslant n\) and that, if X denotes the matrix of u with respect to any two bases of E and F, there exist no scalar \(\mu \neq 0\) such that the products by \(\mu\) of all the minors of X of order m are zero.
\end{enumerate}

\begin{enumerate}
\def\labelenumi{\arabic{enumi}.}
\setcounter{enumi}{16}
\tightlist
\item
  Let
\end{enumerate}

\[
\sum_ {j = 1} ^ {n} \alpha_ {i j} \xi_ {j} = \beta_ {i} \quad (1 \leqslant i \leqslant m)\tag{*}
\]

be a system of m linear equations in n unknowns over a commutative ring A. Let \(x_{j}\) ( \(1 \leqslant j \leqslant n\) ) be the columns of the matrix \(X = (\alpha_{ij})\) of this system and \(y = (\beta_{i})\) ; suppose that in X all the minors of order \textgreater p are zero but that \(x_{1} \wedge x_{2} \wedge \cdots \wedge x_{p} \neq 0\) . For the system (*) to have a solution, it is necessary that \(x_{1} \wedge x_{2} \wedge \cdots \wedge x_{p} \wedge y = 0\) . Conversely, if this condition holds, there exist \(n + 1\) elements \(\xi_{j}\) ( \(1 \leqslant j \leqslant n + 1\) ) of A such that \(\xi_{n+1} \neq 0\) and

\[
\sum_ {j = 1} ^ {n} \alpha_ {i j} \xi_ {j} = \beta_ {i} \xi_ {n + 1} \quad (1 \leqslant i \leqslant m).
\]

\begin{enumerate}
\def\labelenumi{\arabic{enumi}.}
\setcounter{enumi}{17}
\tightlist
\item
  Let E and F be two free A-modules of respective dimensions m and n, \(u:E \rightarrow F\) a linear mapping and X the matrix of u with respect to any two bases of E and F.
\end{enumerate}

\begin{enumerate}
\def\labelenumi{(\alph{enumi})}
\item
  If \(m \geqslant n\) and there exists a minor of \(X\) of order \(n\) which is invertible, \(u\) is surjective.
\item
  Conversely, if \(u\) is surjective, show that \(m \geqslant n\) and that there exists a minor of \(X\) of order \(n\) which is non-zero. If further, in the ring \(A\) , the ideal generated by the non-invertible elements is distinct from \(A\) , there exists a minor of \(X\) of order \(n\) which is invertible (consider the exterior power \(\bigwedge^{n}u\) ).
\item
  Let B be a commutative ring and A the product ring \(B \times B\) . Give an example of a surjective linear mapping of the A-module \(A^{2}\) onto A, all the elements of whose matrix (with respect to the canonical bases of \(A^{2}\) and A) are divisors of zero.
\end{enumerate}

\begin{enumerate}
\def\labelenumi{\arabic{enumi}.}
\setcounter{enumi}{18}
\item
  Let X be a matrix over a commutative field. For X to be of rank p, it suffices that there exist a minor of X of order p which is \(\neq0\) and such that all the minors of order \(p+1\) containing this minor of order p are zero (show that every column of X is then a linear combination of the p columns to which the elements of the minor of order p in question belong).
\item
  Let \(\mathbf{A}\) be a commutative ring, \(\mathbf{M}\) an arbitrary A-module and \(U = (\alpha_{ij})\) a matrix of type \((m, n)\) with elements in \(\mathbf{A}\) .
\end{enumerate}

\begin{enumerate}
\def\labelenumi{(\alph{enumi})}
\item
  Suppose that \(m = n\) and let \(\Delta = \det(U)\) ; then, if \(x_1, \ldots, x_n\) are elements of \(M\) such that \(\sum_{j=1}^{n} \alpha_{ij} x_j = 0\) for \(1 \leqslant i \leqslant n\) , then \(\Delta x_i = 0\) for every index \(i\) .
\item
  Suppose that \(m\) and \(n\) are arbitrary, that all the minors of \(U\) of order \(> r\) are zero and that a minor of order \(r\) is invertible, for example that of the matrix \((\alpha_{ij})\) where \(i \leqslant r\) and \(j \leqslant r\) ; then if \(a_i\) denotes the row of index \(i\) in \(U\) , the \(a_i\) of index \(\leqslant r\) form a basis of the submodule of \(A^n\) generated by all the rows of \(U\) ; for \(i > r\) , therefore \(a_i = \sum_{k=1}^{r} \lambda_{ik} a_k\) . Then let \(y_1, \ldots, y_m\) be elements of \(M\) ; for there to exist elements \(x_1, \ldots, x_n\) of \(M\) satisfying the system of equations \(\sum_{j=1}^{n} \alpha_{ij} x_j = y_i\) for \(1 \leqslant i \leqslant m\) , it is necessary and sufficient that \(y_i = \sum_{k=1}^{r} \lambda_{ik} y_k\) for \(i > r\) .
\end{enumerate}

¶ 21. Let K be an infinite commutative field, m, n, r three integers \(\geqslant0\) such that \(r\leqslant n\leqslant m\) , \(\mathbf{M}_{m,n}(\mathbf{K})\) the vector K-space (of dimension mn) consisting of the matrices of type \((m,n)\) over K and V a vector subspace of \(\mathbf{M}_{m,n}(\mathbf{K})\) such that r is the greatest value of the rank of matrices \(X\in V\) . We propose to prove the inequality

\[
\dim (\mathbf {V}) \leqslant m r.\tag{*}
\]

\begin{enumerate}
\def\labelenumi{(\alph{enumi})}
\tightlist
\item
  Show that we can assume that \(m = n\) and that the matrix
\end{enumerate}

\[
X _ {0} = \left( \begin{array}{c c} I _ {r} & 0 \\ 0 & 0 \end{array} \right)
\]

belongs to V. Deduce that every matrix \(X \in \mathbf{V}\) is then of the form

\[
\left( \begin{array}{c c} X _ {1 1} & X _ {1 2} \\ X _ {2 1} & 0 \end{array} \right)
\]

(where \(X_{11}\) is of order \(r\) ) with the condition \(X_{21}X_{12} = 0\) (use the fact that for all \(\xi \in \mathbf{K}\) , the matrix \(\xi X_0 + X\) is of rank \(\leqslant r\) ).

\begin{enumerate}
\def\labelenumi{(\alph{enumi})}
\setcounter{enumi}{1}
\tightlist
\item
  Deduce from (a) that, for two matrices \(X\) , \(Y\) in \(\mathbf{V}\) ,
\end{enumerate}

\[
X _ {2 1} Y _ {1 2} + Y _ {2 1} X _ {1 2} = 0.
\]

\begin{enumerate}
\def\labelenumi{(\alph{enumi})}
\setcounter{enumi}{2}
\tightlist
\item
  For every matrix \(X \in \mathbf{V}\) , we write \(F(X) = (X_{11}X_{12}) \in \mathbf{M}_{r,n}(\mathbf{K})\) ; \(f\) is a linear mapping and \(\dim(\mathbf{V}) = \dim(\operatorname{Im}(f)) + \dim(\operatorname{Ker}(f))\) . On the other hand, let \(u_x\) denote the linear form \((\mathrm{Y}_{11}\mathrm{Y}_{12}) \mapsto \operatorname{Tr}(\mathrm{X}_{21}\mathrm{Y}_{12})\) on \(\mathbf{M}_{r,n}(\mathbf{K})\) . Show that the linear mapping \(X \mapsto u_X\) of \(\operatorname{Ker}(f)\) into the dual \((\mathbf{M}_{r,n}(\mathbf{K}))^*\) is injective; complete the proof of (*) by noting that the image under \(X \mapsto u_X\) of \(\operatorname{Ker}(f)\) is contained in the orthogonal of \(\operatorname{Im}(f)\) .
\end{enumerate}

¶ 22. Let F be a mapping of \(\mathbf{M}_{n}(\mathrm{K})\) into a set E, where K is a commutative field, such that, for every triple of matrices \(X\) , \(Y\) , \(Z\) in \(\mathbf{M}_{n}(\mathrm{K})\) ,

\[
\mathbf {F} (X Y Z) = \mathbf {F} (X Z Y).
\]

We except the case \(\mathbf{M}_{2}(\mathbf{F}_{2})\) ; show then that there exists a mapping \(\Phi\) of K into E such that \(\mathrm{F}(X)=\Phi(\det(X))\) . (Using the Corollary to Proposition 17 (no. 9), show first that, for \(U\in\mathbf{SL}_{n}(\mathrm{K})\) , \(\mathrm{F}(X)=\mathrm{F}(XU)\) ; deduce that, in \(\mathbf{GL}_{n}(\mathrm{K})\) , \(\mathrm{F}(X)\) depends only on \(\det(X)\) . On the other hand, if, for \(r\leqslant n-1\) , we write

\[
X _ {r} = \left( \begin{array}{c c} I _ {r} & 0 \\ 0 & 0 \end{array} \right)
\]

show that there exists a matrix \(Y\) of rank \(r\) such that \(X_{r-1} = YX_r\) and \(Y = X_rY\) and deduce that \(F(X)\) has the same value for all matrices of rank \(< n\) .

\begin{enumerate}
\def\labelenumi{\arabic{enumi}.}
\setcounter{enumi}{22}
\tightlist
\item
  In every A-algebra E, if \(x_{1}, \ldots, x_{n}\) are elements of E, we write
\end{enumerate}

\[
\left[ x _ {1} x _ {2} \dots x _ {n} \right] = \sum_ {\sigma \in \mathfrak {S} _ {n}} \varepsilon_ {\sigma} x _ {\sigma (1)} x _ {\sigma (2)} \dots x _ {\sigma (n)}.
\]

Show that if E admits a finite generating system over A, there exists an integer N such that \([x_{1}x_{2}\ldots x_{N}] = 0\) for all elements \(x_{j} (1 \leqslant j \leqslant N)\) of E.

\begin{enumerate}
\def\labelenumi{\arabic{enumi}.}
\setcounter{enumi}{23}
\tightlist
\item
  Let \(X_{i}\) ( \(1 \leqslant i \leqslant m\) ) be square matrices of the same order \(n\) over the commutative ring A. Show that if \(m\) is even then (Exercise 23)
\end{enumerate}

\[
\operatorname{Tr} \left[ X _ {1} X _ {2} \dots X _ {m} \right] = 0
\]

and if \(m\) is odd

\[
\operatorname{Tr} \left[ X _ {1} X _ {2} \dots X _ {m} \right] = m. \operatorname{Tr} \left(X _ {m} \left[ X _ {1} X _ {2} \dots X _ {m - 1} \right]\right).
\]

(If \(\lambda\) is the cyclic permutation \((1,2,\ldots,m)\) , C the cyclic subgroup of \(\mathfrak{S}_m\) generated by \(\lambda\) and H the subgroup of \(\mathfrak{S}_m\) consisting of the permutations leaving \(m\) invariant, note that every permutation in \(\mathfrak{S}_m\) can be written uniquely as \(\sigma\tau\) with \(\sigma\in\mathrm{H}\) and \(\tau\in\mathrm{C}\) ; then use the fact that the trace of a product of matrices is invariant under cyclic permutation of the factors.)

¶ 25. Let A be a commutative ring containing the field Q of rational numbers and let X be a square matrix of even order \(2n \geqslant 2\) over A. Let L be the subset \(\{1, 2, \ldots, n + 1\}\) of \(\{1, 2, \ldots, 2n\}\) and L' its complement. Show the following formula (``Redei's identity''):

\[
(*) \quad \det (X _ {\mathrm{L}, \mathrm{L}}) \det (X _ {\mathrm{L} ^ {\prime}, \mathrm{L} ^ {\prime}}) = \sum_ {\mathrm{H}, \mathrm{K}} (- 1) \frac {r - 1}{n \binom {n - 1} {r}} \det (X _ {\mathrm{H}, \mathrm{K}}) \det (X _ {\mathrm{H} ^ {\prime}, \mathrm{K} ^ {\prime}})
\]

where H runs through the set of subsets of n elements of \(\{1,2,\ldots,2n\}\) containing l and K the set of subsets of n elements of L and \(r=\operatorname{Card}(H\cap L')\) . (Evaluate the coefficient of a term \(x_{\sigma(1),1}x_{\sigma(2),2}\ldots x_{\sigma(2n),2n}\) in the right-hand side of (*) for an arbitrary permutation \(\sigma\in S_{2n}\) , noting that this coefficient can only be \(\neq0\) if \(H=\sigma(K)\) ; show that it can only be \(\neq0\) if \(\sigma(L)=L\) .)

\begin{enumerate}
\def\labelenumi{\arabic{enumi}.}
\setcounter{enumi}{25}
\tightlist
\item
  \begin{enumerate}
  \def\labelenumii{(\alph{enumii})}
  \tightlist
  \item
    In the notation of Proposition 19 (no. 10), suppose that there exist two A{[}X{]}-linear mappings \(g_1, g_2\) of M{[}X{]} into M'{[}X{]} such that \(\psi' \circ g_2 = g_1 \circ \psi\) ; show that there then exists an A{[}X{]}-linear mapping \(g\) of \(M_u\) into \(M_{u'}\) such that \(\phi' \circ g_1 = \phi' \circ \bar{g}\) . Further, if \(g_1\) and \(g_2\) are A{[}X{]}-isomorphisms of M{[}X{]} onto M'{[}X{]}, \(g\) is an A{[}X{]}-isomorphism of \(M_u\) onto \(M_{u'}\) .
  \end{enumerate}
\end{enumerate}

\begin{enumerate}
\def\labelenumi{(\alph{enumi})}
\setcounter{enumi}{1}
\tightlist
\item
  In the notation of no. 10, the endomorphisms \(u, u'\) are called equivalent if there exist two isomorphisms \(f_1, f_2\) of M onto M such that \(u' \circ f_2 = f_1 \circ u\) ; they are called similar if there exists an isomorphism \(f\) of M onto M' such that \(u' \circ f = f \circ u\) . Deduce from (a) that, for \(u\) and \(u'\) to be similar, it is necessary and sufficient that the endomorphisms \(X - \bar{u}\) and \(X - \bar{u}'\) (of the A{[}X{]}-modules M{[}X{]} and M'{[}X{]} respectively) be equivalent.
\end{enumerate}

\BourbakiExercisesSection

\begin{enumerate}
\def\labelenumi{\arabic{enumi}.}
\tightlist
\item
  Let K be a commutative field with at least 3 elements and A the algebra over K with a basis consisting of the unit element 1 and two elements \(e_1, e_2\) such that \(e_1^2 = e_1, e_1e_2 = e_2, e_2e_1 = e_2^2 = 0\) ; let \(\mathbf{A}^0\) be the opposite algebra to A. Show that there exists in A elements \(x\) such that \(\operatorname{Tr}_{\mathbf{A} / \mathbf{K}}(x) \neq \operatorname{Tr}_{\mathbf{A}^0 / \mathbf{K}}(x)\) and \(\mathbf{N}_{\mathbf{A} / \mathbf{K}}(x) \neq \mathbf{N}_{\mathbf{A}^0 / \mathbf{K}}(x)\) .
\end{enumerate}

\BourbakiExercisesSection

\begin{enumerate}
\def\labelenumi{\arabic{enumi}.}
\tightlist
\item
  If A and B are K-algebras and \(\alpha\) and \(\beta\) two K-homomorphisms of A into B, an \((\alpha, \beta)\) -derivation of A into B is a K-linear mapping d of A into B satisfying the relation
\end{enumerate}

\[
d (x y) = (d x) \alpha (y) + \beta (x) (d y) \quad \text { for } \quad x, y \quad \text { in } \quad A,
\]

in other words, a derivation in the sense of Definition 1 (no. 2) where \(\mathbf{A} = \mathbf{A}' = \mathbf{A}''\) , \(\mathbf{B} = \mathbf{B}' = \mathbf{B}''\) , \(d = d' = d''\) , \(\mu: \mathbf{A} \times \mathbf{A} \to \mathbf{A}\) is multiplication, \(\lambda_1: \mathbf{B} \times \mathbf{A} \to \mathbf{B}\) is the K-bilinear mapping \((z, x) \mapsto z\alpha(x)\) and \(\lambda_2: \mathbf{A} \times \mathbf{B} \to \mathbf{B}\) the K-bilinear mapping \((x, z) \mapsto \beta(x)z\) .

Suppose henceforth that A and B are commutative fields and that \(\alpha \neq \beta\) . Show that there then exists an element \(b \in \mathbf{B}\) such that \(d\) is the \((\alpha, \beta)\) -derivation \(x \mapsto b\alpha(x) - \beta(x)b\) . (It can be assumed that \(d \neq 0\) . Show first that the kernel N of \(d\) is the set of \(x \in \mathbf{A}\) such that \(\alpha(x) = \beta(x)\) and then show that for \(x \notin \mathbf{N}\) the element \(b = d(x)/( \alpha(x) - \beta(x))\) is independent of the choice of \(x\) .)

\begin{enumerate}
\def\labelenumi{\arabic{enumi}.}
\setcounter{enumi}{1}
\tightlist
\item
  Let K be a commutative field of characteristic \(\neq2\) , A and B two K-algebras and \(\alpha\) a K-linear mapping of A into B; an \(\alpha\) -derivation of A into B is a K-linear mapping d of A into B satisfying the relation
\end{enumerate}

\[
d (x y) = (d x) \alpha (y) + \alpha (x) (d y) \quad \text { for } \quad x, y \quad \text { in } \quad A,
\]

in other words, a derivation on the sense of Definition 1 with \(\mathbf{A} = \mathbf{A}' = \mathbf{A}''\) , \(\mathbf{B} = \mathbf{B}' = \mathbf{B}''\) , \(d = d' = d''\) , \(\mu\) multiplication in \(\mathbf{A}\) , \(\lambda_1\) the mapping \((z,x) \mapsto z\alpha(x)\) of \(\mathbf{B} \times \mathbf{A}\) into \(\mathbf{B}\) and \(\lambda_2\) the mapping \((x,z) \mapsto \alpha(x)z\) of \(\mathbf{A} \times \mathbf{B}\) into \(\mathbf{B}\) .

Show that if B is a commutative field, an extension of K, and \(d \neq 0\) , there exists \(b \neq 0\) in B such that

\[
\alpha (x y) = \alpha (x) \alpha (y) + b (d x) (d y) \quad \text { for } \quad x, y \quad \text { in } \quad A.
\]

If \(b = c^2\) with \(c \in \mathbf{B}\) , \(c \neq 0\) , there exist two homomorphisms \(f, g\) of \(\mathbf{A}\) into \(\mathbf{B}\) such that \(\alpha(x) = (f(x) + g(x)) / 2\) , \(dx = (f(x) - g(x)) / 2c\) . Otherwise, there exists a quadratic algebra (§ 2, no. 3) \(\mathbf{B}'\) over \(\mathbf{B}\) , a \(\mu \in \mathbf{B}'\) such that \(\mu^2 = c\) , and a homomorphism \(f\) of \(\mathbf{A}\) into \(\mathbf{B}'\) such that \(\alpha(x) = (f(x) + \overline{f(x)}) / 2\) , \(dx = (f(x) - \overline{f(x)}) / 2\) .

¶ 3. Let K be a field, commutative or otherwise, E the left vector space \(\mathrm{K}_{s}^{(\mathbf{N})}\) and let \((e_{n})_{n\in\mathbf{N}}\) be the canonical basis of this space. On the other hand let \(\sigma\) be an endomorphism of K and d a \((\sigma,1_{\mathbf{K}})\) -derivation (Exercise 1) of K into itself, in other words an additive mapping of K into itself such that

\[
d (\xi \eta) = (d \xi) \sigma (\eta) + \xi (d \eta).
\]

The elements of E are linear combinations of products \(\alpha_{n}e_{n}\) and a multiplication is defined on E (Z-bilinear mapping of E × E into E) by the conditions

\[
\begin{array}{l} (\alpha e _ {n}) (\beta e _ {m}) = (\alpha e _ {n - 1}) (\sigma (\beta) e _ {m + 1} + (d \beta) e _ {m}) \text { for } n \geqslant 1 \\ (\alpha e _ {0}) (\beta e _ {m}) = (\alpha \beta) e _ {m}. \end{array}
\]

Thus an (in general non-commutative) ring structure is defined on E with \(e_{0}\) as unit element (identified with the element l of K). If we write \(X = e_{1}\) , then \(e_{n} = X^{n}\) for \(n \geqslant 1\) , so that the elements of E can be written as

\[
\alpha_ {0} + \alpha_ {1} X + \dots + \alpha_ {n} X ^ {n},
\]

with the multiplication rule

\[
\mathrm{X} \alpha = \sigma (\alpha) \mathrm{X} + (d \alpha) \quad \text { for } \quad \alpha \in \mathrm{K}.
\]

This ring is denoted by K{[}X; σ, d{]} and is called the ring of non-commutative polynomials in X with coefficients in K, relative to σ and d.~For a non-zero element \(z = \sum_{j} \alpha_{j} X^{j}\) of E, the degree deg(z) is the greatest integer j such that \(\alpha_{j} \neq 0\) .

\begin{enumerate}
\def\labelenumi{(\alph{enumi})}
\tightlist
\item
  Show that if \(u, v\) are two elements \(\neq 0\) of \(E\) , then:
\end{enumerate}

\[
\begin{array}{r l} u v \neq 0 & \text { and } \quad \deg (u v) = \deg (u) + \deg (v); \\ \deg (u + v) & \leqslant \sup (\deg (u), \deg (v)) \quad \text { if } u + v \neq 0. \end{array}
\]

Moreover, if \(\deg(u) \geqslant \deg(v)\) , there exist two elements w, r of E such that \(u = wv + r\) and, either r = 0, or \(\deg(r) < \deg(v)\) (``Euclidean division'' in E).

\begin{enumerate}
\def\labelenumi{(\alph{enumi})}
\setcounter{enumi}{1}
\tightlist
\item
  Conversely, let R be a ring such that there exists a mapping \(u \mapsto \deg(u)\) of the set \(R'\) of elements \(\neq 0\) of R into N, satisfying the conditions of (a). Show that the set K consisting of 0 and the non-zero elements \(u \in R\) such that \(\deg(u) = 0\) is a (not necessarily commutative) field. If \(K \neq R\) and \(x\) is an element of R for which the integer \(\deg(x) > 0\) is the smallest possible, show
\end{enumerate}

that every element of R can be written uniquely in the form \(\sum_{i} \alpha_j x^j\) with \(\alpha_j \in K\) . (To see that every element of R is of this form, argue by reductio ad absurdum by considering an element \(u \in R\) which is not of this form and for which \(\deg(u)\) is the smallest possible.) Prove that there exist an endomorphism \(\sigma\) of K and a \((\sigma, 1_K)\) -derivation \(d\) of K such that \(x\alpha = \sigma(\alpha)x + d\alpha\) for all \(\alpha \in K\) . Deduce that R is isomorphic to a field or a ring K{[}X; \(\sigma, d]\) .

\begin{enumerate}
\def\labelenumi{\arabic{enumi}.}
\setcounter{enumi}{3}
\tightlist
\item
  Let M be a graded K-module of type N; show that every graded endomorphism of degree r can be extended uniquely to a derivation of degree r of the universal anticommutative algebra \(\mathbf{G}(\mathbf{M})\) (§ 7, Exercise 13).
\end{enumerate}

* 5. Let K be a field of characteristic \(p > 0\) , B the field K(X) of rational functions in one indeterminate and A the subfield K(X \(^{p}\) ) of B. Show that the canonical homomorphism \(\Omega_{\mathbf{K}}(\mathbf{A}) \otimes_{\mathbf{A}} \mathbf{B} \to \Omega_{\mathbf{K}}(\mathbf{B})\) is not injective.*

\BourbakiExercisesSection

\begin{enumerate}
\def\labelenumi{\arabic{enumi}.}
\item
  If V is a finitely generated projective A-module, so is End(V), canonically identified with \(V \otimes_{A} V^{*}\) (II, §4, no. 2, Corollary to Proposition 2). If every element \(u \in \text{End}(V)\) is associated with the A-linear form \(v \mapsto \text{Tr}(uv)\) on End(V), an A-linear bijection is defined of End(V) onto its dual (End(V))*. Identifying End(V) and its dual A-module under this mapping, the A-algebra structure on End(V) defines by duality (no. 1, Example 3) an A-cogebra structure on End(V), which is coassociative and has as counit the form Tr: \(V \to A\) . In particular, for \(V = A^{n}\) , a cogebra structure is defined on \(\mathbf{M}_{n}(\mathbf{A})\) for which the coproduct is given by \(c(E_{ij}) = \sum_{k} E_{kj} \otimes E_{ik}\) . This cogebra structure and the algebra structure on \(\mathbf{M}_{n}(\mathbf{A})\) do not define a bigebra structure.
\item
  Show that in Definition 3 (no. 4), condition (3) (relative to graded bigebras) is equivalent to the following: the product \(m: \mathbf{E} \otimes \mathbf{E} \to \mathbf{E}\) is a morphism of the graded cogebra \(\mathbf{E} \otimes \mathbf{E}\) into the graded cogebra \(\mathbf{E}\) .
\item
  Show that if \(\mathbf{M}\) is an A-module there exists on \(\mathsf{T}(\mathbf{M})\) one and only one cocommutative bigebra structure whose algebra structure is the usual structure and whose coproduct \(c\) is such that, for all \(x\in \mathbf{M}\) , \(c(x) = 1\otimes x + x\otimes 1\) .
\item
  Let \(\mathbf{E}\) be a commutative bigebra over \(\mathbf{A}\) , \(m\) the product \(\mathbf{E} \otimes \mathbf{E} \to \mathbf{E}\) , \(c\) the coproduct \(\mathbf{E} \to \mathbf{E} \otimes \mathbf{E}\) , \(e\) the unit element and \(\gamma\) the counit of \(\mathbf{E}\) .
\end{enumerate}

\begin{enumerate}
\def\labelenumi{(\alph{enumi})}
\tightlist
\item
  Show that for every commutative A-algebra B, the law of composition which associates with every ordered pair
\end{enumerate}

\[
(u, v) \in \operatorname{Hom} _ {\mathbf {A} - \mathrm{alg.}} (\mathbf {E}, \mathbf {B}) \times \operatorname{Hom} _ {\mathbf {A} - \mathrm{alg.}} (\mathbf {E}, \mathbf {B})
\]

the A-algebra homomorphism

\[
\mathrm{E} \xrightarrow {c} \mathrm{E} \otimes \mathrm{E} \xrightarrow {u \otimes v} \mathrm{B} \otimes \mathrm{B} \xrightarrow {m _ {\mathrm{B}}} \mathrm{B}
\]

(where \(m_{\mathbf{B}}\) is the product in B) defines a monoid structure on \(\operatorname{Hom}_{\mathbf{A}\text{-alg.}}(\mathbf{E}, \mathbf{B})\) . (b) In order that, for every commutative A-algebra B, \(\operatorname{Hom}_{\mathbf{A}\text{-alg.}}(\mathbf{E}, \mathbf{B})\) be a group, it is necessary and sufficient that there exist an A-algebra homomorphism \(i: \mathbf{E} \to \mathbf{E}\) such that \(i(e) = e\) and the composite mappings \(m \circ (1_{\mathbf{E}} \otimes i) \circ c\) and \(m \circ (i \otimes 1_{\mathbf{E}}) \circ c\) are both equal to the mapping \(x \mapsto \gamma(x)e\) . \(i\) is then called an inversion in E; it is an isomorphism of E onto the opposite bigebra \(E^0\) .

\begin{enumerate}
\def\labelenumi{\arabic{enumi}.}
\setcounter{enumi}{4}
\item
  Let \(E_{1}\) , \(E_{2}\) be two bigebras over A. Show that on the tensor product \(E_{1} \otimes_{A} E_{2}\) the algebra and cogebra structures the tensor product of those on \(E_{1}\) and \(E_{2}\) define on \(E_{1} \otimes_{A} E_{2}\) a bigebra structure.
\item
  Let M be a graded A-module of type N. Define a skew graded bigebra structure on the universal anticommutative algebra G(M) (§ 7, Exercise
\end{enumerate}

13); deduce a canonical graded algebra homomorphism \(\mathbf{G}(\mathbf{M}^{*}) \to \mathbf{G}(\mathbf{M})^{*\mathbf{gr}}\) and notions of inner product between elements of \(\mathbf{G}(\mathbf{M})\) and elements of \(\mathbf{G}(\mathbf{M}^{*})\) .

\begin{enumerate}
\def\labelenumi{\arabic{enumi}.}
\setcounter{enumi}{6}
\tightlist
\item
  With the hypotheses and notation of Proposition 11 (no. 9), prove the formula
\end{enumerate}

\[
(\det g) \phi^ {- 1} = \bigwedge (g) \circ \phi^ {- 1} \circ \bigwedge (^ {t} g).
\]

Form this formula and (79) (no. 11) derive a new proof of the expression for the inverse of a square matrix using the matrix of its cofactors (§ 8, no. 6, formula (26)).

\begin{enumerate}
\def\labelenumi{\arabic{enumi}.}
\setcounter{enumi}{7}
\tightlist
\item
  Let E be a free A-module of dimension n.~For every A-linear mapping v of \(\bigwedge(\mathbf{M})\) into an A-module G, we can write \(v = \sum_{p=0}^{n} v_{p}\) , where \(v_{p}\) is the restriction of v to \(\bigwedge^{p}(\mathbf{M})\) ; we then set \(\eta v = \sum_{p=0}^{n} (-1)^{p} v_{p}\) , so that \(\eta^{2} v = v\) . Let u be an isomorphism of M onto its dual \(M^{*}\) and let \(\Delta\) be the determinant of the matrix of u with respect to a basis \((e_{i})\) of M and the dual basis of \(M^{*}\) (a determinant which does not depend on the basis of M chosen). If \(\phi\) is the isomorphism of \(\bigwedge(\mathbf{M})\) onto \(\bigwedge(\mathbf{M}^{*})\) defined starting with \(e = e_{1} \wedge e_{2} \wedge \cdots \wedge e_{n}\) , show the formula
\end{enumerate}

\[
\eta^ {n + 1} \Delta \phi = \bigwedge (t u) \circ \phi^ {- 1} \circ \bigwedge (u).
\]

\begin{enumerate}
\def\labelenumi{\arabic{enumi}.}
\setcounter{enumi}{8}
\tightlist
\item
  The adjoint of a square matrix \(X\) of order \(n\) over \(A\) is the matrix \(\tilde{X} = (\det(X^{it}))\) of minors of \(X\) of order \(n - 1\) . Show that if \(X\) is invertible then \(\det(\tilde{X}) = (\det X)^{n-1}\) and that every minor \(\det(\tilde{X}_{H,K})\) of \(\tilde{X}\) of order \(p\) is given by the formula
\end{enumerate}

\[
\det (\tilde {X} _ {\mathrm{H}, \mathrm{K}}) = (\det X) ^ {p - 1} (X _ {\mathrm{H} ^ {\prime}, \mathrm{K} ^ {\prime}})\tag{1}
\]

where H' and K' are the complements of H and K respectively in \([1, n]\) (``Jacobi's identities''; use relation (80) of no. 11).

\begin{enumerate}
\def\labelenumi{\arabic{enumi}.}
\setcounter{enumi}{9}
\tightlist
\item
  From every identity \(\Phi = 0\) between minors of an arbitrary invertible square matrix \(X\) of order \(n\) over A another identity \(\tilde{\Phi} = 0\) can be derived called the complement of \(\Phi = 0\) , by applying the identity \(\Phi = 0\) to the minors of the adjoint \(\tilde{X}\) of \(X\) and then replacing the minors of \(\tilde{X}\) by functions of the minors of \(X\) using identity (1) of Exercise 9. In this way show the following identity:
\end{enumerate}

\[
\det (X ^ {i h}) \det (X ^ {j k}) - \det (X ^ {i k}) \det (X ^ {j h}) = (\det X) \det (X ^ {i j, h k}) \quad \text { for } \quad i <   j
\]

where \(X^{ij, hk}\) denotes the minor of X of order n - 2 obtained by suppressing in X the rows of index i, j and the columns of index h, k.

¶ 11. Let \(\Phi = 0\) be an identity between minors of an invertible square matrix \(X\) of order \(n\) over a commutative ring A, \(\tilde{\Phi} = 0\) the complementary identity (Exercise 10) and \(k\) an integer \(>0\) . Let \(Y\) be an invertible square matrix of order \(n + k\) and let \(\tilde{Y}_0\) be the submatrix of the adjoint matrix \(\tilde{Y}\) formed by suppressing in \(\tilde{Y}\) the rows of index \(\leqslant k\) and the columns of index \(\leqslant k\) . If \(\tilde{Y}_0\) is assumed to be invertible and the identity \(\tilde{\Phi} = 0\) is applied to the minors of \(\tilde{Y}_0\) and then each minor which appears in this identity (considered as a minor of \(\tilde{Y}\) ) is replaced by its expression as a function of the minors of \(Y\) , using identity (1) of Exercise 9, an identity \(\Phi_k = 0\) is obtained between minors of the matrix \(Y\) (valid if \(Y\) and \(\tilde{Y}_0\) are invertible) which is called the extension of order \(k\) of the identity \(\Phi = 0\) .

In particular, Let \(A = (\alpha_{ij})\) be an invertible square matrix of order \(n + k\) , B the submatrix of A of order k obtained by suppressing in A the rows and columns of index \textgreater k and \(\Delta_{ij}\) the determinant of the matrix of order \(k + 1\) obtained by suppressing in A the rows of index \textgreater k except that of index \(k + i\) and the columns of index \textgreater k except that of index \(k + j\) ; if C denotes the matrix \((\Delta_{ij})\) of order n, show the identity

\[
\det C = (\det A) (\det B) ^ {n - 1}
\]

valid when \(A\) and \(B\) are invertible (show that it is the extension of order \(k\) of the total expansion of a determinant of order \(n\) ).

State the identities obtained by extending the Laplace expansion (§ 8, no. 6) and the identity of Exercise 2 of § 8.

¶ 12. Let \(X = (\xi_{ij})\) be a square matrix of order \(n\) over a commutative field \(K\) ; let \(H\) be a subset of \([1, n]\) with \(p\) elements and \(H'\) the complement of \(H\) in \([1, n]\) ; suppose that, for every ordered pair of indices \(h \in H\) , \(k \in H'\) ,

\[
\sum_ {i} \xi_ {i h} \xi_ {i k} = 0.
\]

Show that, for all subsets L, M of \([1, n]\) with p elements,

\[
\rho_ {\mathbf {M}, \mathbf {M} ^ {\prime}} \det (X _ {\mathrm{L}, \mathrm{H}}) \det (X _ {\mathbf {M} ^ {\prime}, \mathbf {H} ^ {\prime}}) - \rho_ {\mathrm{L}, \mathrm{L} ^ {\prime}} \det (X _ {\mathbf {M}, \mathbf {H}}) \det (X _ {\mathrm{L}, \mathrm{H} ^ {\prime}}) = 0.
\]

(Consider the columns \(x_h\) of index \(h \in \mathbf{H}\) in \(X\) as vectors in \(\mathbf{E} = \mathbf{K}^n\) and the columns \(x_k^*\) of index \(k \in \mathbf{H}'\) as vectors in \(\mathbf{E}^*\) and show that the \((n - p)\) -form \(x_{\mathbf{H}'}^*\) is proportional to \(\phi_p(x_{\mathbf{H}})\) .)

\begin{enumerate}
\def\labelenumi{\arabic{enumi}.}
\setcounter{enumi}{12}
\tightlist
\item
  Let \(\Gamma\) and \(\Delta\) be two determinants of invertible matrices of order n over a commutative ring and \(\Gamma_{ij}\) the determinant obtained by replacing in \(\Gamma\) the column of index i by the column of \(\Delta\) of index j. Show that
\end{enumerate}

\[
\det (\Gamma_ {i j}) = \Gamma^ {n - 1} \Delta .
\]

(Expand \(\Gamma_{ij}\) by the \(i\) -th column and use Exercise 9.)

\begin{enumerate}
\def\labelenumi{\arabic{enumi}.}
\setcounter{enumi}{13}
\item
  Let M be a free A-module of dimension n \textgreater{} 1. Show that every isomorphism \(\psi\) of \(\bigwedge^{p}(M)\) onto \(\bigwedge^{n-p}(M^{*})\) such that \(\bigwedge^{n-p}(\check{u}) \circ \psi = \psi \circ \bigwedge^{p}(u)\) for every automorphism u of M, of determinant equal to 1, is one of the isomorphisms \(\phi_{p}\) defined in Proposition 12 of no. 11 (proceed as in §7, Exercise 4).
\item
  Let E be a free A-module of dimension n.~Let z be a p-vector over E and \(z^{*}\) a \((p + q)\) -form on E; z can be canonically identified (§ 7, Exercise 8) with the skewsymmetrization \(az_{1}\) of a contravariant tensor \(z_{1}\) of order p and \(z^{*}\) with the skewsymmetrization of a covariant tensor of order \(p + q\) . If, in the mixed tensor \(z_{1}\) . \(z^{*}\) , the k-th contravariant index and the \((p + k)\) -th covariant index are contracted for \(1 \leqslant k \leqslant p\) , show that the covariant tensor of order q thus obtained is skewsymmetrized and is canonically identified with the q-form \(z^{*} \perp z\) , up to sign.
\item
  Let E be a free A-module of dimension n.~The regressive product of \(x \in \bigwedge(\mathrm{E})\) and \(y \in \bigwedge(\mathrm{E})\) with respect to a basis \(\{e\}\) of \(\bigwedge^{n}(\mathrm{E})\) , denoted by \(x \vee y\) , is the element \(\phi(\phi(x) \wedge \phi(y))\) , the isomorphism \(\phi\) being relative to e. This product is only defined to within an invertible factor, depending on the basis \(\{e\}\) chosen. Show that if x is a p-vector and y a q-vector, \(x \vee y = 0\) if \(p + q < n\) and that, if \(p + q \geqslant n\) , \(x \vee y\) is a \((p + q - n)\) -vector such that
\end{enumerate}

\[
y \vee x = (- 1) ^ {(n - p) (n - q)} x \vee y.
\]

The regressive product is associative and distributive with respect to addition and defines on \(\bigwedge(\mathrm{E})\) a unital algebra structure isomorphic to the exterior algebra structure. Express the components of \(x \vee y\) as a function of those of \(x\) and \(y\) for a given basis on \(\mathbf{E}\) .

\begin{enumerate}
\def\labelenumi{\arabic{enumi}.}
\setcounter{enumi}{16}
\tightlist
\item
  Let K be a commutative field and E a vector space over K.
\end{enumerate}

\begin{enumerate}
\def\labelenumi{(\alph{enumi})}
\item
  Let \(z\) be a \(p\) -vector \(\neq 0\) over \(E\) and let \(V_{z}\) be the vector subspace of \(E\) consisting of the vectors \(x\) such that \(z \wedge x = 0\) . Then \(\dim(V_{z}) \leqslant p\) ; if \((x_{i})_{1 \leqslant i \leqslant q}\) is a free system of vectors of \(V_{z}\) , there exists a \((p - q)\) -vector \(v\) such that \(z = v \wedge x_{1} \wedge \dots \wedge x_{q}\) . If \(E\) is finite-dimensional, then \(V_{z} = \phi(M_{z}^{0})\) , where \(M_{z}^{0}\) is the orthogonal in \(E^{*}\) of the subspace \(M_{z}\) associated with \(z\) . For \(z\) to be a pure \(p\) -vector, it is necessary and sufficient that \(\dim(V_{z}) = p\) and then \(V_{z} = M_{z}\) .
\item
  Let \(v\) be a pure \(p\) -vector, \(w\) a pure \(q\) -vector and \(V\) and \(W\) the subspaces of \(E\) associated respectively with \(v\) and \(w\) . In order that \(V \subset W\) , it is necessary and sufficient that \(q \geqslant p\) and that there exist a \((q - p)\) -vector \(u\) such that \(w = u \wedge v\) . In order that \(V \cap W = \{0\}\) , it is necessary and sufficient that \(v \wedge w \neq 0\) ; the subspace \(V + W\) is then associated with the pure \((p + q)\) -vector \(v \wedge w\) .
\item
  Let \(u\) and \(v\) be two pure \(p\) -vectors and \(U\) and \(V\) the subspaces associated respectively with \(u\) and \(v\) . For \(u + v\) to be pure, it is necessary and sufficient that \(\dim(U \cap V) \geqslant p - 1\) .
\end{enumerate}

\begin{enumerate}
\def\labelenumi{\arabic{enumi}.}
\setcounter{enumi}{17}
\item
  Let \(z = \sum_{H} \alpha_{H} e_{H}\) be a non-zero pure p-vector of an n-dimensional vector space E, expressed using its components with respect to an arbitrary basis \((e_{i})_{1 \leqslant i \leqslant n}\) of E. Let G be a subset of \([1, n]\) with p elements such that \(\alpha_{G} \neq 0\) ; let \((i_{h})_{1 \leqslant h \leqslant p}\) be the sequence of indices in G arranged in increasing order and \((j_{k})_{1 \leqslant k \leqslant n - p}\) the sequence of indices in the complement \(G'\) of G, arranged in increasing order. For every ordered pair \((h, k)\) of indices such that \(1 \leqslant h \leqslant p\) , \(1 \leqslant k \leqslant n - p\) , let \(\beta_{hk}\) be the component \(\alpha_{H}\) of z corresponding to the subset H of \([1, n]\) consisting of the p - 1 indices in G distinct from \(i_{h}\) and of index \(j_{k} \in G'\) ; let X be the matrix \((\beta_{hk})\) with p rows and n - p columns. Given an arbitrary subset L of \([1, n]\) with p elements such that \(L \cap G\) has \(q \geqslant 1\) elements, show that \((\alpha_{G})^{q-1} \alpha_{L}\) is equal, up to sign, to the minor of X of order q, consisting of the rows of index h such that \(i_{h} \in G \cap G\) and the columns of index k such that \(j_{k} \in L \cap G\) . (Write z as being of the form \(\alpha_{G} y_{1} \wedge y_{2} \wedge \cdots \wedge y_{p}\) , where the vectors \(y_{i}\) are such that, in the matrix Y with n rows and p columns whose columns are the \(y_{i}\) , the submatrix consisting of the rows whose indices belong to G is the unit matrix.)
\item
  Over a finite-dimensional vector space E, let \(z = x_{1} \wedge x_{2} \wedge \cdots \wedge x_{p}\) be a pure p-vector \(\neq 0\) ; let \(v^{*}\) be a q-form ( \(q \leqslant p\) ) and \(u^{*}\) an arbitrary element of \(\bigwedge(\mathrm{E}^{*})\) . Show that
\end{enumerate}

\[
(u ^ {*} \wedge v ^ {*}) \lrcorner z = (- 1) ^ {q (q - 1) / 2} \sum_ {\mathrm{H}} \rho_ {\mathrm{H}, \mathrm{K}} \langle v ^ {*}, x _ {\mathrm{H}} \rangle (u ^ {*} \lrcorner x _ {\mathrm{K}})
\]

where H runs through the set of subsets of \([1, p]\) with q elements, K is the complement of H in \([1, p]\) and \(x_{H}\) (resp. \(x_{K}\) ) denotes the exterior product of the \(x_{i}\) such that \(i \in H\) (resp. \(i \in K\) ).

\begin{enumerate}
\def\labelenumi{\arabic{enumi}.}
\setcounter{enumi}{19}
\item
  Let E be an n-dimensional vector space, z a pure p-vector over E, V the vector subspace of E associated with z, \(u^{*}\) a pure q-vector over \(E^{*}\) , \(W'\) the subspace of \(E^{*}\) associated with \(u^{*}\) and W the subspace of E, of dimension n - q, orthogonal to \(W'\) . If q \textless{} p, \(u^{*} \lrcorner z\) is a pure (p - q)-vector over E, which is zero if \(\dim(V \cap W) > p - q\) ; if \(\dim(V \cap W) = p - q\) , \(V \cap W\) is associated with \(u^{*} \lrcorner z\) .
\item
  \begin{enumerate}
  \def\labelenumii{(\alph{enumii})}
  \tightlist
  \item
    Show directly (without using the isomorphisms \(\phi_p\) ) that every \((n - 1)\) -vector over an \(n\) -dimensional vector space is pure.
  \end{enumerate}
\end{enumerate}

\begin{enumerate}
\def\labelenumi{(\alph{enumi})}
\setcounter{enumi}{1}
\tightlist
\item
  Let A be a commutative algebra over a field K with a basis consisting of the unit element 1 and three elements \(a_{1}\) , \(a_{2}\) , \(a_{3}\) whose products with one another are zero. Let E be the A-module \(\mathbf{A}^3\) and \((e_i)_{1\leqslant i\leqslant 3}\) its canonical basis; show that the bivector
\end{enumerate}

\[
a _ {1} (e _ {2} \land e _ {3}) + a _ {2} (e _ {3} \land e _ {1}) + a _ {3} (e _ {1} \land e _ {2})
\]

is not pure.

\begin{enumerate}
\def\labelenumi{\arabic{enumi}.}
\setcounter{enumi}{21}
\tightlist
\item
  Let E be an ordered set in which every interval \([a, b]\) is finite. Let S be the subset of \(E \times E\) consisting of the ordered pairs \((x, y)\) such that \(x \leqslant y\) . Show that a coassociative cogebra structure is defined on the A-module \(\mathbf{A}^{(\mathrm{s})}\) of formal linear combinations of the elements of S by taking as coproduct
\end{enumerate}

\[
c ((x, y)) = \sum_ {x \leqslant z \leqslant y} (x, z) \otimes (z, y).
\]

For this cogebra the linear form \(\gamma\) such that

\[
\gamma ((x, y)) = 0 \quad \text { if } \quad x \neq y \text { in   E }, \quad \gamma ((x, x)) = 1
\]

is a counit.

\begin{enumerate}
\def\labelenumi{\arabic{enumi}.}
\setcounter{enumi}{22}
\tightlist
\item
  Let C be a cogebra over a ring A. A sub-A-module S of C is called a subcogebra of C if \(c(S) \subset \operatorname{Im}(S \otimes S)\) .
\end{enumerate}

\begin{enumerate}
\def\labelenumi{(\alph{enumi})}
\item
  If C is coassociative (resp. cocommutative), so is every subcogebra. If \(\gamma\) is a counit of C, the restriction of \(\gamma\) to any subcogebra of C is a counit.
\item
  Which are the subcogebras of the cogebra \(\mathbf{A}^{(\mathbf{X})}\) (no. 1, Example 4).?
\item
  If \((\mathbf{S}_{\alpha})\) is a family of subcogebras of \(\mathbf{C}\) , \(\sum_{\alpha} \mathbf{S}_{\alpha}\) is a subcogebra of \(\mathbf{C}\) (the smallest containing all the \(\mathbf{S}_{\alpha}\) ).
\end{enumerate}

\begin{enumerate}
\def\labelenumi{\arabic{enumi}.}
\setcounter{enumi}{23}
\tightlist
\item
  Let C be a cogebra over a ring A. A sub-A-module S of C is called a right (resp. left) coideal if \(c(S) \subset \operatorname{Im}(S \otimes C)\) (resp. \(c(S) \subset \operatorname{Im}(C \otimes S)\) ).
\end{enumerate}

\begin{enumerate}
\def\labelenumi{(\alph{enumi})}
\item
  A subcogebra (Exercise 23) is a right and left coideal. The converse is true if \(\mathbf{A}\) is a field.
\item
  Every sum of right (resp. left) coideals is a right (resp. left) coideal.
\item
  If S is a right (resp. left) coideal of C, the orthogonal \(S^{0}\) is a right (resp. left) ideal of the dual algebra \(C^{*}\) .
\item
  If A is a field and \(\mathfrak{I}'\) a right (resp. left) ideal of the dual algebra \(\mathbf{C}^*\) , the orthogonal \(\mathfrak{I}''^0\) in \(\mathbf{C}\) is a right (resp. left) coideal of \(\mathbf{C}\) . (Argue by reductio ad absurdum by assuming that there exists an \(x \in \mathfrak{I}''^0\) such that \(c(x) \notin \mathfrak{I}''^0 \otimes \mathbf{C}\) ; write \(c(x) = \sum_{i} y_i \otimes z_i\) , where the \(z_i\) are linearly independent, and show that there would be \(y' \in \mathfrak{I}'\) and \(z' \in \mathbf{C}^*\) such that \(\langle y' z', x \rangle \neq 0\) .)
\end{enumerate}

Deduce that for a vector subspace S of C to be a right (resp. left) coideal, it is necessary and sufficient that \(S^{0}\) be a right (resp. left) ideal of \(C^{*}\) .

\begin{enumerate}
\def\labelenumi{(\alph{enumi})}
\setcounter{enumi}{4}
\tightlist
\item
  Deduce from (d) that if A is a field every intersection of right coideals (resp. left coideals, resp. subcogebras) of C is a right coideal (resp. left coideal, resp. subcogebra). For S to be a subcogebra of C, it is necessary and sufficient that its orthogonal \(S^{0}\) be a two-sided ideal of the algebra \(C^{*}\) .
\end{enumerate}

\begin{enumerate}
\def\labelenumi{\arabic{enumi}.}
\setcounter{enumi}{24}
\tightlist
\item
  Let C be a cogebra over a ring A. A sub-A-module S of C is called a coideal if \(c(\mathbf{S}) \subset \operatorname{Im}(\mathbf{S} \otimes \mathbf{C}) + \operatorname{Im}(\mathbf{C} \otimes \mathbf{S})\) . If C is counital, with counit \(\gamma\) , a coideal S is called conull if \(\gamma(x) = 0\) for \(x \in S\) .
\end{enumerate}

\begin{enumerate}
\def\labelenumi{(\alph{enumi})}
\item
  A right or left coideal is a coideal. Every sum of coideals is a coideal.
\item
  If S is a coideal of C, the orthogonal \(S^{0}\) is a subalgebra of the dual algebra \(C^{*}\) . If C is counital and S is conull, \(S^{0}\) is unital.
\item
  Suppose that A is a field and C a counital cogebra. Show that if \(E'\) is a unital subalgebra of \(C^{*}\) with the same unit element as \(C^{*}\) , the orthogonal \(E'^{0}\) in C is a conull coideal.
\item
  If S is a coideal of C, show that there exists on C/S one and only one cogebra structure such that the canonical mapping \(p: C \rightarrow C/S\) is a cogebra morphism. If C is counital and S is conull, C/S is counital.
\item
  For every cogebra morphism \(f\colon\mathbf{C}\to\mathbf{C}^{\prime},\mathbf{S}=\operatorname{Ker}(f)\) is a coideal of \(\mathbf{C}\) , \(\mathbf{S}^{\prime}=\operatorname{Im}(f)\) a subcogebra of \(\mathbf{C}^{\prime}\) and, in the canonical factorization
\end{enumerate}

\[
\mathrm{C} \xrightarrow {p} \mathrm{C} / \mathrm{S} \xrightarrow {g} \mathrm{S} ^ {\prime} \xrightarrow {j} \mathrm{C} ^ {\prime}
\]

of \(f, g\) is a cogebra isomorphism.

\begin{enumerate}
\def\labelenumi{(\arabic{enumi})}
\setcounter{enumi}{25}
\tightlist
\item
  Let C be a coassociative counital cogebra over a commutative field K.
\end{enumerate}

\begin{enumerate}
\def\labelenumi{(\alph{enumi})}
\item
  For every vector subspace E of C which is a left C*-module under the law \((u, x) \mapsto u \perp x\) , the left annihilator of this C*-module is a two-sided ideal \(S'\) of the algebra \(C^*\) . Show that if E is finite-dimensional over K, \(C^*/S'\) is finite-dimensional.
\item
  Deduce from (a) that for every element \(x \in \mathbb{C}\) , the subcogebra of \(\mathbb{C}\) generated by \(x\) is finite-dimensional. (Note that the \(\mathbb{C}^*\) -module \(\mathbb{E}\) generated by \(x\) is finite-dimensional and that, if \(\mathfrak{S}'\) is its left annihilator, \(x \in \mathfrak{S}'^0\) , the orthogonal of \(\mathfrak{S}'\) in \(\mathbb{C}\) .)
\end{enumerate}

\begin{enumerate}
\def\labelenumi{(\arabic{enumi})}
\setcounter{enumi}{26}
\tightlist
\item
  Let B be an associative unital algebra over a commutative field K and let \(m: B \otimes_{K} B \to B\) be the K-linear mapping defining the multiplication on B. The tensor product \(B^{*} \otimes_{K} B^{*}\) (where \(B^{*}\) is the dual vector space of the vector space B) is canonically identified with a vector subspace of \((\mathbf{B} \otimes_{\mathbf{K}} \mathbf{B})^{*}\) (II, §7, no. 7, Proposition 16).
\end{enumerate}

\begin{enumerate}
\def\labelenumi{(\alph{enumi})}
\item
  Show that for an element \(w \in \mathbf{B}^*\) the following conditions are equivalent: (1) \({}^t m(w) \in \mathbf{B}^* \otimes \mathbf{B}^*\) ; (2) the set of \(x \lrcorner w\) , where \(x\) runs through \(\mathbf{B}\) , is a finite-dimensional subspace of \(\mathbf{B}^*\) ; (3) the set of \(w \lfloor x\) , where \(x\) runs through \(\mathbf{B}\) , is a finite-dimensional subspace of \(\mathbf{B}^*\) ; (4) the set of \(x \lrcorner w \lfloor y\) , where \(x\) and \(y\) run through \(\mathbf{B}\) , is contained in a finite-dimensional subspace of \(\mathbf{B}^*\) .
\item
  Let \(\mathbf{B}' \subset \mathbf{B}^*\) be the set of \(w \in \mathbf{B}^*\) satisfying the equivalent conditions in
\end{enumerate}

(a). Show that \({}^t m(\mathbf{B}') \subset \mathbf{B}' \otimes_{\mathbb{K}} \mathbf{B}'\) . (For \(w \in \mathbf{B}'\) , write \({}^t m(w) = \sum_i u_i \otimes v_i\) , where for example the \(u_i\) are linearly independent in \(\mathbf{B}^*\) ; show that then \(v_i \in \mathbf{B}'\) for all \(i\) using the associativity of \(\mathbf{B}\) and arguing by reductio ad absurdum.) \(\mathbf{B}'\) is therefore the largest cogebra contained in \(\mathbf{B}^*\) , for which \({}^t m\) is the coproduct. \(\mathbf{B}'\) is called the dual cogebra of the algebra \(\mathbf{B}\) . When \(\mathbf{B}\) is finite-dimensional, \(\mathbf{B}' = \mathbf{B}^*\) .

\begin{enumerate}
\def\labelenumi{(\alph{enumi})}
\setcounter{enumi}{2}
\item
  Show that if \(\mathbf{B}\) is a commutative field of infinite dimension over \(\mathbf{K}\) , \(\mathbf{B}' = \{0\}\) . (Note that the orthogonal of the set of \(x \perp w\) , for same \(w \in \mathbf{B}^*\) , \(x\) running through \(\mathbf{B}\) , is an ideal of \(\mathbf{B}\) .)
\item
  Let \(\mathbf{B}_1, \mathbf{B}_2\) be two K-algebras and \(f: \mathbf{B}_1 \to \mathbf{B}_2\) a K-homomorphism of algebras. Show that \(^t f(\mathbf{B}_2') \subset \mathbf{B}_1'\) and that \(^t f\) , restricted to \(\mathbf{B}_2'\) , is a cogebra homomorphism of \(\mathbf{B}_2'\) into \(\mathbf{B}_1'\) .
\item
  If \(\mathbf{C}\) is a coassociative counital cogebra over \(\mathbf{K}\) , the canonical injection \(\mathbf{C} \to \mathbf{C}^{**}\) of the vector space \(\mathbf{C}\) into its bidual maps \(\mathbf{C}\) onto a subspace of \((\mathbf{C}^*)'\) , the dual cogebra of the algebra \(\mathbf{C}^*\) dual to \(\mathbf{C}\) , and is a cogebra homomorphism of \(\mathbf{C}\) into \((\mathbf{C}^*)'\) .
\end{enumerate}

\BourbakiAppendixExercises

\begin{enumerate}
\def\labelenumi{\arabic{enumi}.}
\item
  Show that in alternative Cayley A-algebra (notation of § 2, no. 4) \(\mathrm{T}((xy)z) = \mathrm{T}(x(yz))\) . (Reduce it to proving that \(a(x,y,z) = a(\bar{z},\bar{y},\bar{x})\) .)
\item
  Let F be an alternative A-algebra, in which the relation 2x = 0 implies x = 0. Show that, for all x, y, z in F, \(x(yz)x = (xy)(zx)\) (Moufang's identity). (In the identity
\end{enumerate}

\[
(x y) z + y (z x) = x (y z) + (y z) x
\]

successively replace \(y\) by \(xy\) and \(z\) by \(zx\) .

\begin{enumerate}
\def\labelenumi{\arabic{enumi}.}
\setcounter{enumi}{2}
\tightlist
\item
  Let \(\mathbf{F}\) be an alternative Cayley A-algebra, in which the relation \(2x = 0\) implies \(x = 0\) . Show that if \(x \in \mathbf{F}\) is invertible then, for all \(y, z, t\) in \(\mathbf{F}\) ,
\end{enumerate}

\[
(\mathrm{N} (x) + \mathrm{N} (y)) (\mathrm{N} (z) + \mathrm{N} (t)) = \mathrm{N} (x \bar {z} + y \bar {t}) + \mathrm{N} (x t - (x z) (x ^ {- 1} y))
\]

(use Exercises 1 and 2).
